% =====================================================================
%  Testing, Experimentation and Validation
%  A SECTION of the FYP report. Include from the parent chapter with:
%      % =====================================================================
%  Testing, Experimentation and Validation  (~3 pages)
%  A SECTION of the FYP report. Include from the parent chapter with:
%      % =====================================================================
%  Testing, Experimentation and Validation  (~3 pages)
%  A SECTION of the FYP report. Include from the parent chapter with:
%      % =====================================================================
%  Testing, Experimentation and Validation  (~3 pages)
%  A SECTION of the FYP report. Include from the parent chapter with:
%      \input{testing_validation}
%
%  Intended to follow \input{hitl_edge_processing}; cross-references
%  \ref{sec:hitl}. A longer version with the full 26-case test matrix is
%  kept in testing_validation_long.tex if an appendix is wanted.
%
%  REQUIRED PACKAGES
%      \usepackage{booktabs}
%      \usepackage{array}
%      \usepackage{graphicx}
%      \usepackage{subcaption}
%      \graphicspath{{figures/}}
%
%  IMAGES EXPECTED in docs/report/figures/  (see checklist at end)
%      detection_a.png  detection_b.png  detection_c.png
% =====================================================================

\section{Testing, Experimentation and Validation}
\label{sec:test}

Testing of the hardware-in-the-loop testbed described in
Section~\ref{sec:hitl} was organised around a difficulty specific to this
kind of system: a fault at a low layer does not raise an error at the layer
above it, but produces plausible and wrong behaviour there. A subscriber
starved of buffer space receives nothing and reports nothing; two nodes with
mismatched topic names run indefinitely without complaint. Testing was
therefore performed bottom-up --- physical link, path separation, middleware
transport, detector, end-to-end --- with each layer required to pass before
the next was attempted.

Three results are reported here. The first establishes that the two-link
architecture does what it claims. The second establishes that the detector
finds real people rather than plausible-looking texture. The third concerns
the measurement chain itself, and is included because it determines whether
the other two can be believed.

% ---------------------------------------------------------------------
\subsection{Test environment}
\label{sec:test-env}

Both machines run Ubuntu~22.04 with ROS\,2 Humble and Fast~DDS~2.6, so the
middleware is identical at each end of the link and cannot itself be a
source of difference. The simulation host runs Gazebo Classic in lockstep
with ArduPilot SITL, and ns-3 with TapBridge for the wireless channel. The
edge node is a Raspberry~Pi~4B running YOLOv8n restricted to the COCO
\emph{person} class. The camera produces $1280 \times 720$ frames of
$2.76$\,MB at $5$\,Hz. The test scene contains exactly five human models at
coordinates fixed in the world definition, which makes recall a counted
quantity rather than an estimated one.

The Pi's single Ethernet port is divided by 802.1Q tagging into
\texttt{eth0.10} ($10.0.0.2$, camera, unimpaired) and \texttt{eth0.42}
($10.42.0.12$, detections, through ns-3), created by
\texttt{scripts/netns/pi\_hitl\_link.sh}. The matching host interfaces
\texttt{eth-cam} and \texttt{eth-rf} are created by step~1c of
\texttt{scripts/netns/launch\_single\_uav\_netns.sh}, the latter bridged
into \texttt{tap-uav2}.

% ---------------------------------------------------------------------
\subsection{Result 1: the two links behave differently}
\label{sec:test-links}

The central claim of the architecture is that detection results cross a
simulated wireless channel while camera imagery does not. Two paths that
both work are not evidence of this; they must also be shown to differ in the
intended way, and imagery must be shown to be absent from the radio path.

Table~\ref{tab:test-links} reports sixty ICMP samples on the sensor path and
thirty on the radio path. Large samples are necessary: the first packet on the
radio path includes address resolution across the simulated channel, and a
two-sample measurement of the same configuration returned a mean of
$167$\,ms against $74.5$\,ms for thirty. The short sample is biased, not
merely noisy.

\begin{table}[htbp]
\centering
\caption{Measured characteristics of the two paths. Sensor path over sixty
samples, radio path over thirty; zero packet loss on both.}
\label{tab:test-links}
\begin{tabular}{@{}lrrrr@{}}
\toprule
Path & Min (ms) & Mean (ms) & Max (ms) & Jitter (ms) \\
\midrule
Sensor, VLAN 10 (wired)   & 1.0  & 1.4  & 1.8   & 0.12 \\
Radio, VLAN 42 (via ns-3) & 10.1 & 74.5 & 158.0 & 38.4 \\
\midrule
Ratio                     & --   & $53\times$ & -- & $320\times$ \\
\bottomrule
\end{tabular}
\end{table}

The radio path is approximately \(53\times\) slower on average, and its
variation is larger by more than two orders of magnitude.

Its internal spread is itself informative: a minimum of \(10.1\)\,ms against
a mean of \(74.5\)\,ms is a range of some \(7\times\) on one link, where a
channel applying a fixed delay would produce a narrow distribution. The
variation indicates that ns-3 is modelling contention, in which packets queue
and back off, rather than adding a constant offset.

The radio measurement was repeated. An independent set of thirty samples taken
several hours earlier, with the host reconfigured in between, gave a mean of
\(77.4\)\,ms against the \(74.5\)\,ms reported here, agreeing within
\(4\%\).

The sensor figures were taken \emph{while} the camera was streaming at
\(110\)\,Mbps. An idle-cable measurement over the same number of samples
gives the same mean, so the two logical links sharing one physical cable do
not delay one another. Queue management was verified rather than assumed:
both machines run \texttt{fq\_codel} with a \(5\)\,ms target, and byte
queue limits on the edge node had reduced its driver ring to approximately one
packet.

Absence of imagery on the radio path was established negatively, by
observation rather than by inspecting a rule set, since separation here is
enforced by addressing and bridge membership rather than by filtering. No
$10.0.0.0/24$ traffic was observed on \texttt{tap-uav2}, the point at which
traffic enters ns-3, and the byte counters on the two interfaces differ by
approximately four orders of magnitude --- consistent with $2.76$\,MB per
frame on one path against $118$\,B per detection message on the other, a
reduction of approximately \(23{,}000\times\).

A structural argument supports the measurements. The ground station receiver
executes inside the \texttt{gcsns} network namespace, whose only interface
is one end of a veth pair attached to \texttt{tap-gcs}. That namespace has
no interface on $10.0.0.0/24$ and no route to it. A message arriving at the
receiver cannot have taken any other path, because no other path exists from
inside the namespace. Separation is therefore a property of the topology
rather than of a setting that could be silently overridden.

% ---------------------------------------------------------------------
\subsection{Result 2: detections correspond to real subjects}
\label{sec:test-detection}

Visual confirmation was made a precondition of accepting any quantitative
result. A detector reporting a plausible count while placing its boxes on
building texture would satisfy every numerical check in this section, and an
earlier configuration did exactly that: sixteen confident detections against
five subjects, none correct. Only inspection of annotated frames
distinguishes the two cases. Those false positives were eliminated by
raising the confidence threshold and narrowing the camera field of view.

\begin{figure}[htbp]
\centering
\begin{subfigure}[b]{0.32\textwidth}
  \includegraphics[width=\textwidth]{detection_a.png}
  \caption{Two subjects, confidence $0.59$ and $0.42$.}
  \label{fig:test-det-a}
\end{subfigure}
\hfill
\begin{subfigure}[b]{0.32\textwidth}
  \includegraphics[width=\textwidth]{detection_b.png}
  \caption{One subject, confidence $0.48$.}
  \label{fig:test-det-b}
\end{subfigure}
\hfill
\begin{subfigure}[b]{0.32\textwidth}
  \includegraphics[width=\textwidth]{detection_c.png}
  \caption{Three subjects, confidence $0.59$, $0.47$ and $0.46$.}
  \label{fig:test-det-c}
\end{subfigure}
\caption{Detector output during flight, with the edge node's log above each
view. Bounding boxes coincide with the human models rather than with ground
texture. The log records every frame, including those with no detection, so
that a stalled node is distinguishable from an empty scene, and reports
inference duration per frame.}
\label{fig:test-detection}
\end{figure}

Figure~\ref{fig:test-detection} shows three moments from a patrol flight.
Between one and four of the five subjects were detected per frame, at
confidences of \(0.42\)--\(0.59\). The shortfall against the known count of
five is attributable to scene geometry rather than to detector deficiency:
the models stand approximately one metre apart, so from the oblique camera
view the nearer figures partially occlude those behind them, and a partially
occluded subject at approximately \(53\) pixels in height falls below the
confidence threshold. Because the number of subjects is known exactly, any
count exceeding five would be a false positive without the need to
adjudicate individual boxes; none was observed in this configuration.

Inference duration was measured around the model invocation alone, excluding
message reception, format conversion and publication, so that the figure
remains comparable between back-ends and against standalone execution. Over
\(480\) consecutive frames of one patrol flight the mean was
\(1{,}093\)\,ms, with a minimum of \(1{,}025\)\,ms. The distribution is tight
apart from a single outlier at \(6{,}949\)\,ms, one frame in \(480\),
attributable to contention for the processor from another task rather than to
the model itself.

The Pi therefore completes rather less than one detection per second against a
camera delivering five frames per second, and discards approximately \(80\%\)
of what it receives --- which is why the camera subscription uses
\texttt{BEST\_EFFORT} with queue depth one, keeping the newest frame rather
than accumulating a backlog of stale ones.

% ---------------------------------------------------------------------
\subsection{Result 3: validating the measurement chain}
\label{sec:test-chain}

The two results above depend on the camera actually reaching the edge node
at the rate the configuration specifies. That assumption failed, and the
manner of its failure is reported here because it determines what the other
measurements are worth.

The symptom was that the edge node received imagery at \(0.19\)\,Hz where
the camera was configured for \(5\)\,Hz, a shortfall of approximately
\(26\times\). The visible consequences appeared three layers above the
cause: the detector reported no humans in \(93\) of \(94\) frames, and the
patrol mission timed out at every waypoint. Both invite explanation in terms
of detector settings or flight control, and neither has anything to do with
either --- the aircraft flew correctly and the detector functioned
correctly, but the scene was being observed through roughly one frame every
five seconds.

\begin{table}[htbp]
\centering
\caption{Eliminating candidate causes of the frame-rate shortfall. Each row
is a measurement that excludes one layer.}
\label{tab:test-elimination}
\begin{tabular}{@{}p{38mm}p{38mm}p{40mm}@{}}
\toprule
Hypothesis & Measurement & Outcome \\
\midrule
Simulation running slowly
  & Real-time factor from \texttt{/clock}
  & \(0.998\); simulated time tracks wall clock \\[3pt]
Gazebo rendering slowly
  & Host-side subscriber on the camera topic
  & \(5.00\)\,Hz, gaps \(0.19\)--\(0.21\)\,s; source correct \\[3pt]
Link negotiated at $100$\,Mbps
  & \texttt{ethtool} on the edge node
  & $1000$\,Mb/s full duplex \\[3pt]
Fragments lost in flight
  & Interface counters, both machines
  & Zero errors, zero drops \\[3pt]
Receive buffers too small
  & \texttt{net.core.rmem\_max}
  & $512$\,MB, already raised \\[3pt]
\textbf{Send buffers too small}
  & \texttt{net.core.wmem\_max}
  & \(\mathbf{208}\)\,\textbf{kB}, left at the default \\
\bottomrule
\end{tabular}
\end{table}

The DDS profile requests a \(16\)\,MB send buffer. The kernel caps any such
request at \texttt{net.core.wmem\_max}, and performs the reduction silently.
A single frame is therefore approximately \(13\times\) larger than the
buffer available to transmit it: the publisher filled that buffer
immediately and discarded the remaining fragments at the socket layer,
before they reached the network interface. Comparing production with
transmission confirmed this --- the host generated \(414\)\,MB of imagery
over thirty seconds and placed \(10\)\,MB of it on the wire, \(2\)\,Mbps
against the \(110\)\,Mbps the configuration implies. Raising the parameter
on both machines and restarting the stack, which is required because
Fast~DDS applies socket options only at participant creation, restored
delivery to \(5.000\)\,Hz with an inter-arrival standard deviation of
\(4.5\)\,ms: an exact match to the source.

Two properties of this fault are worth recording. First, no counter
registers it. Loss at the socket layer is invisible to interface
statistics, and \texttt{tx\_dropped} on the sending interface and
\texttt{rx\_errors} on the receiving one both remained at zero throughout,
because the discarded fragments were never transmitted. A diagnosis founded
on interface counters would have concluded that the network was healthy,
which it was. Second, the two directions are governed by separate kernel
parameters. The receive path had been corrected during earlier work and the
transmit path had not, because the failure that prompted the original
correction pointed only at the receiving side. The pre-flight checks in
\texttt{scripts/netns/run\_hitl.sh} now verify both parameters on both
machines before the stack starts.

% ---------------------------------------------------------------------
\subsection{Limitations}
\label{sec:test-limitations}

Clock synchronisation between the two machines has not been established.
End-to-end latency would be computed from timestamps taken on different
hosts, and without synchronisation such values are not meaningful while
remaining plausible in magnitude. End-to-end latency is therefore not
reported, rather than reported with a caveat.

The edge configuration is characterised; the corresponding ground
configuration, in which compressed imagery crosses the simulated channel and
inference is performed at the ground station, remains to be measured. All
measurements were taken with the aircraft close to the ground station,
giving zero packet loss on both paths, so the reliability policy applied to
detection messages is untested under the loss conditions it exists to
handle. The testbed supports a single physical edge node; multi-node
operation has not been evaluated.

% =====================================================================
%  IMAGE CHECKLIST — save into docs/report/figures/
%
%  detection_a.png / detection_b.png / detection_c.png
%      The three screenshots already captured: the detector log above and
%      the rqt_image_view annotated frame below, showing 2, 1 and 3
%      detections respectively.
%
%      CROP EACH ONE before including. At three-to-a-row the full desktop
%      screenshot will be illegible. Keep only:
%        - the last ~8 log lines (enough to show the count and the
%          inference time), and
%        - the image pane itself, without the rqt toolbar and topic
%          selector.
%      Crop all three to the same aspect ratio so the row lines up.
%
%      If the log text is still too small at 0.32\textwidth, drop the log
%      from the figure entirely and show only the three annotated views —
%      the inference figures are already stated in the prose.
%
%  OPTIONAL, if a fourth image is wanted:
%      testbed_hardware.jpg — a photograph of the laptop, the USB gigabit
%      adapter, the single Ethernet cable and the Pi 4B. This is the only
%      image that shows the work is hardware-in-the-loop rather than pure
%      simulation. It would go in Section \ref{sec:test-env}.
% =====================================================================

%
%  Intended to follow % =====================================================================
%  Hardware-in-the-Loop Edge Processing
%  A SECTION of the FYP report. Include from the parent chapter with:
%      \input{hitl_edge_processing}
%
%  HEADING LEVELS USED
%      \section     - the HITL section itself
%      \subsection  - its parts
%      \paragraph   - short findings within Results
%
%  REQUIRED PACKAGES (add to the main preamble if not already present)
%      \usepackage{booktabs}
%      \usepackage{graphicx}
%      \usepackage{tikz}
%      \usetikzlibrary{arrows.meta, positioning, fit, backgrounds}
%
%  All labels are prefixed hitl / tab:hitl- / fig:hitl- so they cannot
%  collide with labels elsewhere in the shared report.
% =====================================================================

\section{Hardware-in-the-Loop Edge Processing}
\label{sec:hitl}

The vision workload was moved off the simulation host and onto physical
hardware, so that the cost of onboard detection could be measured on a processor
representative of one that would actually fly. A Raspberry~Pi~4B performs
detection on imagery produced by the simulation, and returns its results across a
wireless channel simulated in ns-3.

% ---------------------------------------------------------------------
\subsection{Extending the testbed with physical hardware}
\label{sec:hitl-testbed}

The existing stack ran entirely on one machine: Gazebo for the environment and
sensors, ArduPilot SITL for the autopilot, ns-3 for the wireless channel, and
Linux network namespaces isolating the ground station from each aircraft. The
ns-3 scenario defined four nodes --- a ground control station and three aircraft
--- of which the second and third were unused placeholders.

A Raspberry~Pi~4B was introduced as the physical edge node and assigned to the
second aircraft slot. Its TAP device, previously created only to satisfy the
scenario's argument list, was brought into service as the attachment point. The
autopilot, physics and channel simulation remained on the host; only the vision
computation moved onto hardware.

The Pi was configured with Ubuntu~22.04 and ROS\,2 Humble to match the host
exactly, giving both machines the same DDS implementation. The host has no
Ethernet port, so a USB~3.0 gigabit adapter was fitted. An initial low-cost
adapter was replaced after \texttt{iperf3} measured \(4.5\)\,Mbps with \(24\%\)
loss, despite \texttt{ethtool} reporting \(100\)\,Mbps full duplex; the
replacement sustained \(834\)\,Mbps with no retransmissions.

% ---------------------------------------------------------------------
\subsection{Constructing the two-link network}
\label{sec:hitl-network}

Two logically separate paths were required between host and Pi. One carries
camera imagery and must remain unimpaired, since on a real aircraft it is a short
cable inside a single airframe rather than a radio. The other carries detection
results and must pass through the simulated channel, as that is the link the
experiment exists to characterise. The same reasoning applies to the
flight-dynamics traffic between SITL and Gazebo, which also bypasses ns-3.

The Pi provides a single Ethernet interface, so both paths were carried over one
cable using IEEE~802.1Q VLAN tagging. On the host, two virtual interfaces were
created on the physical adapter: \texttt{eth-cam}, tagged VLAN~10 and addressed
\(10.0.0.1\); and \texttt{eth-rf}, tagged VLAN~42, left without an address and
made a member of a new bridge, \texttt{br-uav2}. That bridge also received
\texttt{tap-uav2}, the ns-3 interface for the second aircraft node. On the Pi,
matching interfaces \texttt{eth0.10} (\(10.0.0.2\)) and \texttt{eth0.42}
(\(10.42.0.12\)) were created.

Attaching the physical interface to the bridge directly was rejected during
design. A Linux bridge floods frames to all member ports, so camera traffic would
have been delivered into the ns-3 TAP device as well, injecting approximately
\(111\)\,Mbps of imagery into the simulated channel.

The resulting arrangement enforces separation through routing rather than
filtering. Gazebo is addressed only on \(10.0.0.1\) and can reach the Pi only at
\(10.0.0.2\); the ground station namespace is addressed only on \(10.42.0.10\)
and can reach it only at \(10.42.0.12\). Neither can use the other's path.

\begin{figure}[htbp]
\centering
\begin{tikzpicture}[
  font=\small,
  box/.style={draw, rounded corners=2pt, align=center, inner sep=5pt},
  proc/.style={box, fill=black!4},
  net/.style={box, fill=black!10},
  arr/.style={-{Latex[length=2mm]}, thick},
]
\node[proc] (gz)   {Gazebo\\camera + physics};
\node[proc, below=14mm of gz] (gcs) {Ground station\\(namespace)\\\texttt{10.42.0.10}};
\node[net,  right=16mm of gcs] (ns3) {ns-3\\channel};
\node[net,  above=8mm of ns3]  (br)  {\texttt{br-uav2}\\+ TAP};
\node[proc, right=46mm of gz] (pi) {Raspberry Pi 4B\\YOLOv8n detector};

\draw[arr] (gz) -- node[above, align=center, font=\scriptsize]
  {VLAN 10 \quad camera, 2.76\,MB/frame\\unimpaired} (pi);
\draw[arr] (pi.south) |- node[pos=0.75, above, font=\scriptsize]
  {VLAN 42 \quad detections, 118\,B} (br.east);
\draw[arr] (br) -- (ns3);
\draw[arr] (ns3) -- (gcs);

\begin{scope}[on background layer]
  \node[draw, dashed, rounded corners, fit=(gz)(gcs)(ns3)(br),
        inner sep=7pt, label={[font=\scriptsize\itshape]above left:Host PC}] {};
\end{scope}
\end{tikzpicture}
\caption{The two-link arrangement. Camera imagery reaches the edge node over an
unimpaired wired path; detection results traverse the simulated channel. Both
share one physical cable, separated by VLAN tagging.}
\label{fig:hitl-architecture}
\end{figure}

Addressing was made persistent using NetworkManager on the host and netplan on
the Pi, the two distributions using different network managers. Addresses
assigned imperatively were found to be removed within minutes, as each manager
reconciles kernel state against its own configuration and deletes entries it did
not create.

% ---------------------------------------------------------------------
\subsection{Configuring the middleware}
\label{sec:hitl-dds}

A Fast DDS profile was written and applied on both machines. It pins participants
to the intended interfaces through an interface whitelist, since both machines
are attached to WiFi as well as the wired link and DDS otherwise advertises on
every available interface. The profile also disables the shared-memory transport,
which is not isolated by network namespaces and would allow participants to
exchange data without traversing the simulated channel.

Kernel socket buffers were raised on both machines from the \(208\)\,kB default
to \(512\)\,MB. A single \(1280 \times 720\) frame is \(2.76\)\,MB and fragments
into approximately \(1{,}900\) UDP datagrams, of which the default buffer holds
about \(7\%\); reassembly never completed and subscribers received nothing, with
no error reported.

The camera subscription on the Pi was changed from the ROS\,2 default of
\texttt{RELIABLE} with queue depth ten to \texttt{BEST\_EFFORT} with depth one.
Under the default policy the publisher retains each sample until the subscriber
acknowledges it, and because the Pi acknowledges at approximately \(1\)\,Hz, the
camera output of the entire simulation was throttled to \(0.27\)\,Hz.
\texttt{RELIABLE} was retained for detection messages, which must not be lost.

% ---------------------------------------------------------------------
\subsection{Deploying the detector}
\label{sec:hitl-detector}

The detection node runs YOLOv8n restricted to the COCO \emph{person} class,
subscribing to the camera topic and publishing results as JSON. It was deployed
to the Pi and instrumented to record inference duration, measured around the
model invocation alone so that figures remain comparable between back-ends and
against standalone execution.

Two dependency conflicts were resolved during deployment. Installing
\texttt{ultralytics} brings in NumPy~2.x, against which \texttt{cv\_bridge} --- a
compiled extension built for the 1.x C ABI --- fails. Installing an older OpenCV
to compensate reintroduces NumPy~2 as a dependency, and OpenCV~5.0 separately
breaks \texttt{cv\_bridge} by renaming the type constants it uses to build its
conversion table. Versions were pinned to \texttt{numpy==1.26.4} and
\texttt{opencv-python==4.10.0.84}, the only combination satisfying both.

The model was additionally exported to NCNN for ARM-optimised inference. Because
NCNN fixes the input tensor shape at export time, the export geometry was
specified explicitly as \(544 \times 960\) to match the letterboxed frame; a
default export produces a square input in which approximately \(43\%\) of the
computation is spent on padding, at \(19.9\)\,GFLOPS against \(11.3\)\,GFLOPS for
the correct shape. Half-precision export was avoided, the Cortex-A72 lacking fp16
SIMD support.

% ---------------------------------------------------------------------
\subsection{Automation}
\label{sec:hitl-automation}

Bring-up was consolidated into a single script that starts the namespaces, ns-3,
Gazebo, SITL, the micro-ROS agent and the ground station receiver, creates the
VLANs and bridge, and launches the detector on the Pi over SSH. Pre-flight checks
were added for each failure encountered during development: presence of the host
address before Gazebo starts, since Fast DDS enumerates interfaces only at
participant creation and cannot use one assigned later; reachability of the Pi;
adequate socket buffer limits on both machines; and confirmation after startup
that Gazebo bound the intended interface.

% ---------------------------------------------------------------------
\subsection{Results}
\label{sec:hitl-results}

Detections were verified visually, by inspecting annotated frames, before any
quantitative result was accepted. An earlier configuration had reported up to
sixteen detections against five subjects; these were false positives on building
texture, eliminated by raising the confidence threshold and narrowing the camera
field of view.

\paragraph{Channel characterisation.}
Table~\ref{tab:hitl-links} compares the two paths, each measured over thirty ICMP
samples. Short samples are unreliable: the first packet includes address
resolution across the simulated channel, and a two-sample measurement of the same
configuration returned a mean of \(167\)\,ms against \(59\)\,ms for thirty.

\begin{table}[htbp]
\centering
\caption{Measured characteristics of the two links, thirty samples each.}
\label{tab:hitl-links}
\begin{tabular}{lrrrr}
\toprule
Path & Min (ms) & Mean (ms) & Max (ms) & Jitter (ms) \\
\midrule
Sensor, VLAN 10 (wired)   & 0.4  & 1.5  & 2.7   & 0.4  \\
Radio, VLAN 42 (via ns-3) & 11.7 & 77.4 & 242.6 & 54.8 \\
\bottomrule
\end{tabular}
\end{table}

Packet loss was zero on both paths, the aircraft remaining close to the ground
station throughout. The radio path is approximately \(53\times\) slower than the
wired path on average and varies approximately \(140\times\) more.

The distribution within the radio path is itself informative. Its minimum of
\(11.7\)\,ms against a mean of \(77.4\)\,ms is a spread of some \(6.6\times\) on
one link, where a channel applying a fixed delay would produce a narrow one.
The variation indicates that ns-3 is modelling contention, in which packets
queue and back off before transmission, rather than simply adding a constant
offset to each packet.

\paragraph{Sensor transport.}
With the camera at \(20\)\,Hz the wired path carried \(481\)\,Mbps at
\(42{,}544\) packets per second with zero dropped packets, occupying
approximately \(58\%\) of measured link capacity.

\paragraph{Inference latency.}
Table~\ref{tab:hitl-inference} reports latency for each combination of back-end
and input resolution. Changing back-end at constant resolution gives
approximately \(2.5\times\); halving the pixel count at constant back-end gives
approximately \(2.3\times\). Latency therefore scales close to linearly with
pixel count, a pixel ratio of \(0.47\) producing a time ratio of \(0.43\).

\begin{table}[htbp]
\centering
\caption{Inference latency on the Raspberry Pi 4B.}
\label{tab:hitl-inference}
\begin{tabular}{llrr}
\toprule
Back-end & Input & Latency (ms) & Rate (fps) \\
\midrule
PyTorch & $960 \times 544$ & 2540 & 0.39 \\
PyTorch & $640 \times 384$ & 1100 & 0.91 \\
NCNN    & $960 \times 544$ & 1000 & 1.00 \\
NCNN    & $640 \times 384$ & \multicolumn{2}{l}{not measured (est.\ 470)} \\
\bottomrule
\end{tabular}
\end{table}

\paragraph{Detection performance.}
Between two and four of the five subjects were detected per frame, a recall of
\(40\)--\(80\%\). The shortfall is attributable to scene geometry rather than
detector deficiency: the models are separated by approximately one metre, so from
the oblique camera view the nearer figures partially occlude those behind them,
and partially occluded subjects at approximately \(53\) pixels in height fall
below the confidence threshold. Because the number and position of subjects are
known exactly from the world definition, recall is a directly measurable quantity
rather than an estimate.

\paragraph{Transmitted volume.}
Each detection message is approximately \(118\)\,B, against \(2.76\)\,MB for the
corresponding frame, a reduction of approximately \(23{,}000\times\) in data
crossing the radio.

% ---------------------------------------------------------------------
\subsection{Contention between sensor transport and inference}
\label{sec:hitl-contention}

The two workloads on the edge node proved not to be independent. The same model,
on the same input, required \(623\)\,ms executed on a static image but
\(1{,}343\)\,ms within the live pipeline: a difference of \(720\)\,ms, or
approximately half the available inference throughput.

The edge node received frames at \(19.8\)\,Hz while completing inference at
approximately \(0.74\)\,Hz, discarding roughly \(96\%\) of what it received.
Discarded frames are not free: each requires reassembly of some \(2{,}150\) UDP
fragments, deserialisation of \(2.76\)\,MB and a further \(2.76\)\,MB
colour-space conversion, all competing with inference for the same four
processor cores.

Reducing the camera rate at source from \(20\)\,Hz to \(5\)\,Hz --- still
approximately seven times the rate the detector can consume --- reduced inference
latency to approximately \(1{,}000\)\,ms, recovering about a quarter of the
deficit, and reduced link occupancy from \(481\)\,Mbps to \(111\)\,Mbps. The
application-level frame-rate control provided by the relay node cannot achieve
this, as it discards frames after reception, by which point the transport and
deserialisation cost has already been incurred.

This interaction is not observable in pure simulation, where sensor data is
delivered through shared memory at negligible cost, nor in an isolated benchmark
of the detector, since the interference originates outside it. It appears only
when a physical processor performs both reception and inference concurrently.

% ---------------------------------------------------------------------
\subsection{Limitations}
\label{sec:hitl-limitations}

Clock synchronisation between the two machines has not yet been established.
End-to-end latency is computed from timestamps taken on different hosts, and
without synchronisation those values are not meaningful while remaining plausible
in magnitude; synchronisation is required before end-to-end latency is reported.

The edge configuration is characterised; the corresponding ground configuration,
in which compressed imagery crosses the simulated channel and inference is
performed at the ground station, remains to be measured.

Measurements were taken with the aircraft close to the ground station, giving
zero packet loss throughout. Behaviour as range increases and buildings obstruct
the path has not been characterised. The testbed supports a single physical edge
node; multi-node operation has not been evaluated.
; cross-references
%  \ref{sec:hitl}. A longer version with the full 26-case test matrix is
%  kept in testing_validation_long.tex if an appendix is wanted.
%
%  REQUIRED PACKAGES
%      \usepackage{booktabs}
%      \usepackage{array}
%      \usepackage{graphicx}
%      \usepackage{subcaption}
%      \graphicspath{{figures/}}
%
%  IMAGES EXPECTED in docs/report/figures/  (see checklist at end)
%      detection_a.png  detection_b.png  detection_c.png
% =====================================================================

\section{Testing, Experimentation and Validation}
\label{sec:test}

Testing of the hardware-in-the-loop testbed described in
Section~\ref{sec:hitl} was organised around a difficulty specific to this
kind of system: a fault at a low layer does not raise an error at the layer
above it, but produces plausible and wrong behaviour there. A subscriber
starved of buffer space receives nothing and reports nothing; two nodes with
mismatched topic names run indefinitely without complaint. Testing was
therefore performed bottom-up --- physical link, path separation, middleware
transport, detector, end-to-end --- with each layer required to pass before
the next was attempted.

Three results are reported here. The first establishes that the two-link
architecture does what it claims. The second establishes that the detector
finds real people rather than plausible-looking texture. The third concerns
the measurement chain itself, and is included because it determines whether
the other two can be believed.

% ---------------------------------------------------------------------
\subsection{Test environment}
\label{sec:test-env}

Both machines run Ubuntu~22.04 with ROS\,2 Humble and Fast~DDS~2.6, so the
middleware is identical at each end of the link and cannot itself be a
source of difference. The simulation host runs Gazebo Classic in lockstep
with ArduPilot SITL, and ns-3 with TapBridge for the wireless channel. The
edge node is a Raspberry~Pi~4B running YOLOv8n restricted to the COCO
\emph{person} class. The camera produces $1280 \times 720$ frames of
$2.76$\,MB at $5$\,Hz. The test scene contains exactly five human models at
coordinates fixed in the world definition, which makes recall a counted
quantity rather than an estimated one.

The Pi's single Ethernet port is divided by 802.1Q tagging into
\texttt{eth0.10} ($10.0.0.2$, camera, unimpaired) and \texttt{eth0.42}
($10.42.0.12$, detections, through ns-3), created by
\texttt{scripts/netns/pi\_hitl\_link.sh}. The matching host interfaces
\texttt{eth-cam} and \texttt{eth-rf} are created by step~1c of
\texttt{scripts/netns/launch\_single\_uav\_netns.sh}, the latter bridged
into \texttt{tap-uav2}.

% ---------------------------------------------------------------------
\subsection{Result 1: the two links behave differently}
\label{sec:test-links}

The central claim of the architecture is that detection results cross a
simulated wireless channel while camera imagery does not. Two paths that
both work are not evidence of this; they must also be shown to differ in the
intended way, and imagery must be shown to be absent from the radio path.

Table~\ref{tab:test-links} reports sixty ICMP samples on the sensor path and
thirty on the radio path. Large samples are necessary: the first packet on the
radio path includes address resolution across the simulated channel, and a
two-sample measurement of the same configuration returned a mean of
$167$\,ms against $74.5$\,ms for thirty. The short sample is biased, not
merely noisy.

\begin{table}[htbp]
\centering
\caption{Measured characteristics of the two paths. Sensor path over sixty
samples, radio path over thirty; zero packet loss on both.}
\label{tab:test-links}
\begin{tabular}{@{}lrrrr@{}}
\toprule
Path & Min (ms) & Mean (ms) & Max (ms) & Jitter (ms) \\
\midrule
Sensor, VLAN 10 (wired)   & 1.0  & 1.4  & 1.8   & 0.12 \\
Radio, VLAN 42 (via ns-3) & 10.1 & 74.5 & 158.0 & 38.4 \\
\midrule
Ratio                     & --   & $53\times$ & -- & $320\times$ \\
\bottomrule
\end{tabular}
\end{table}

The radio path is approximately \(53\times\) slower on average, and its
variation is larger by more than two orders of magnitude.

Its internal spread is itself informative: a minimum of \(10.1\)\,ms against
a mean of \(74.5\)\,ms is a range of some \(7\times\) on one link, where a
channel applying a fixed delay would produce a narrow distribution. The
variation indicates that ns-3 is modelling contention, in which packets queue
and back off, rather than adding a constant offset.

The radio measurement was repeated. An independent set of thirty samples taken
several hours earlier, with the host reconfigured in between, gave a mean of
\(77.4\)\,ms against the \(74.5\)\,ms reported here, agreeing within
\(4\%\).

The sensor figures were taken \emph{while} the camera was streaming at
\(110\)\,Mbps. An idle-cable measurement over the same number of samples
gives the same mean, so the two logical links sharing one physical cable do
not delay one another. Queue management was verified rather than assumed:
both machines run \texttt{fq\_codel} with a \(5\)\,ms target, and byte
queue limits on the edge node had reduced its driver ring to approximately one
packet.

Absence of imagery on the radio path was established negatively, by
observation rather than by inspecting a rule set, since separation here is
enforced by addressing and bridge membership rather than by filtering. No
$10.0.0.0/24$ traffic was observed on \texttt{tap-uav2}, the point at which
traffic enters ns-3, and the byte counters on the two interfaces differ by
approximately four orders of magnitude --- consistent with $2.76$\,MB per
frame on one path against $118$\,B per detection message on the other, a
reduction of approximately \(23{,}000\times\).

A structural argument supports the measurements. The ground station receiver
executes inside the \texttt{gcsns} network namespace, whose only interface
is one end of a veth pair attached to \texttt{tap-gcs}. That namespace has
no interface on $10.0.0.0/24$ and no route to it. A message arriving at the
receiver cannot have taken any other path, because no other path exists from
inside the namespace. Separation is therefore a property of the topology
rather than of a setting that could be silently overridden.

% ---------------------------------------------------------------------
\subsection{Result 2: detections correspond to real subjects}
\label{sec:test-detection}

Visual confirmation was made a precondition of accepting any quantitative
result. A detector reporting a plausible count while placing its boxes on
building texture would satisfy every numerical check in this section, and an
earlier configuration did exactly that: sixteen confident detections against
five subjects, none correct. Only inspection of annotated frames
distinguishes the two cases. Those false positives were eliminated by
raising the confidence threshold and narrowing the camera field of view.

\begin{figure}[htbp]
\centering
\begin{subfigure}[b]{0.32\textwidth}
  \includegraphics[width=\textwidth]{detection_a.png}
  \caption{Two subjects, confidence $0.59$ and $0.42$.}
  \label{fig:test-det-a}
\end{subfigure}
\hfill
\begin{subfigure}[b]{0.32\textwidth}
  \includegraphics[width=\textwidth]{detection_b.png}
  \caption{One subject, confidence $0.48$.}
  \label{fig:test-det-b}
\end{subfigure}
\hfill
\begin{subfigure}[b]{0.32\textwidth}
  \includegraphics[width=\textwidth]{detection_c.png}
  \caption{Three subjects, confidence $0.59$, $0.47$ and $0.46$.}
  \label{fig:test-det-c}
\end{subfigure}
\caption{Detector output during flight, with the edge node's log above each
view. Bounding boxes coincide with the human models rather than with ground
texture. The log records every frame, including those with no detection, so
that a stalled node is distinguishable from an empty scene, and reports
inference duration per frame.}
\label{fig:test-detection}
\end{figure}

Figure~\ref{fig:test-detection} shows three moments from a patrol flight.
Between one and four of the five subjects were detected per frame, at
confidences of \(0.42\)--\(0.59\). The shortfall against the known count of
five is attributable to scene geometry rather than to detector deficiency:
the models stand approximately one metre apart, so from the oblique camera
view the nearer figures partially occlude those behind them, and a partially
occluded subject at approximately \(53\) pixels in height falls below the
confidence threshold. Because the number of subjects is known exactly, any
count exceeding five would be a false positive without the need to
adjudicate individual boxes; none was observed in this configuration.

Inference duration was measured around the model invocation alone, excluding
message reception, format conversion and publication, so that the figure
remains comparable between back-ends and against standalone execution. Over
\(480\) consecutive frames of one patrol flight the mean was
\(1{,}093\)\,ms, with a minimum of \(1{,}025\)\,ms. The distribution is tight
apart from a single outlier at \(6{,}949\)\,ms, one frame in \(480\),
attributable to contention for the processor from another task rather than to
the model itself.

The Pi therefore completes rather less than one detection per second against a
camera delivering five frames per second, and discards approximately \(80\%\)
of what it receives --- which is why the camera subscription uses
\texttt{BEST\_EFFORT} with queue depth one, keeping the newest frame rather
than accumulating a backlog of stale ones.

% ---------------------------------------------------------------------
\subsection{Result 3: validating the measurement chain}
\label{sec:test-chain}

The two results above depend on the camera actually reaching the edge node
at the rate the configuration specifies. That assumption failed, and the
manner of its failure is reported here because it determines what the other
measurements are worth.

The symptom was that the edge node received imagery at \(0.19\)\,Hz where
the camera was configured for \(5\)\,Hz, a shortfall of approximately
\(26\times\). The visible consequences appeared three layers above the
cause: the detector reported no humans in \(93\) of \(94\) frames, and the
patrol mission timed out at every waypoint. Both invite explanation in terms
of detector settings or flight control, and neither has anything to do with
either --- the aircraft flew correctly and the detector functioned
correctly, but the scene was being observed through roughly one frame every
five seconds.

\begin{table}[htbp]
\centering
\caption{Eliminating candidate causes of the frame-rate shortfall. Each row
is a measurement that excludes one layer.}
\label{tab:test-elimination}
\begin{tabular}{@{}p{38mm}p{38mm}p{40mm}@{}}
\toprule
Hypothesis & Measurement & Outcome \\
\midrule
Simulation running slowly
  & Real-time factor from \texttt{/clock}
  & \(0.998\); simulated time tracks wall clock \\[3pt]
Gazebo rendering slowly
  & Host-side subscriber on the camera topic
  & \(5.00\)\,Hz, gaps \(0.19\)--\(0.21\)\,s; source correct \\[3pt]
Link negotiated at $100$\,Mbps
  & \texttt{ethtool} on the edge node
  & $1000$\,Mb/s full duplex \\[3pt]
Fragments lost in flight
  & Interface counters, both machines
  & Zero errors, zero drops \\[3pt]
Receive buffers too small
  & \texttt{net.core.rmem\_max}
  & $512$\,MB, already raised \\[3pt]
\textbf{Send buffers too small}
  & \texttt{net.core.wmem\_max}
  & \(\mathbf{208}\)\,\textbf{kB}, left at the default \\
\bottomrule
\end{tabular}
\end{table}

The DDS profile requests a \(16\)\,MB send buffer. The kernel caps any such
request at \texttt{net.core.wmem\_max}, and performs the reduction silently.
A single frame is therefore approximately \(13\times\) larger than the
buffer available to transmit it: the publisher filled that buffer
immediately and discarded the remaining fragments at the socket layer,
before they reached the network interface. Comparing production with
transmission confirmed this --- the host generated \(414\)\,MB of imagery
over thirty seconds and placed \(10\)\,MB of it on the wire, \(2\)\,Mbps
against the \(110\)\,Mbps the configuration implies. Raising the parameter
on both machines and restarting the stack, which is required because
Fast~DDS applies socket options only at participant creation, restored
delivery to \(5.000\)\,Hz with an inter-arrival standard deviation of
\(4.5\)\,ms: an exact match to the source.

Two properties of this fault are worth recording. First, no counter
registers it. Loss at the socket layer is invisible to interface
statistics, and \texttt{tx\_dropped} on the sending interface and
\texttt{rx\_errors} on the receiving one both remained at zero throughout,
because the discarded fragments were never transmitted. A diagnosis founded
on interface counters would have concluded that the network was healthy,
which it was. Second, the two directions are governed by separate kernel
parameters. The receive path had been corrected during earlier work and the
transmit path had not, because the failure that prompted the original
correction pointed only at the receiving side. The pre-flight checks in
\texttt{scripts/netns/run\_hitl.sh} now verify both parameters on both
machines before the stack starts.

% ---------------------------------------------------------------------
\subsection{Limitations}
\label{sec:test-limitations}

Clock synchronisation between the two machines has not been established.
End-to-end latency would be computed from timestamps taken on different
hosts, and without synchronisation such values are not meaningful while
remaining plausible in magnitude. End-to-end latency is therefore not
reported, rather than reported with a caveat.

The edge configuration is characterised; the corresponding ground
configuration, in which compressed imagery crosses the simulated channel and
inference is performed at the ground station, remains to be measured. All
measurements were taken with the aircraft close to the ground station,
giving zero packet loss on both paths, so the reliability policy applied to
detection messages is untested under the loss conditions it exists to
handle. The testbed supports a single physical edge node; multi-node
operation has not been evaluated.

% =====================================================================
%  IMAGE CHECKLIST — save into docs/report/figures/
%
%  detection_a.png / detection_b.png / detection_c.png
%      The three screenshots already captured: the detector log above and
%      the rqt_image_view annotated frame below, showing 2, 1 and 3
%      detections respectively.
%
%      CROP EACH ONE before including. At three-to-a-row the full desktop
%      screenshot will be illegible. Keep only:
%        - the last ~8 log lines (enough to show the count and the
%          inference time), and
%        - the image pane itself, without the rqt toolbar and topic
%          selector.
%      Crop all three to the same aspect ratio so the row lines up.
%
%      If the log text is still too small at 0.32\textwidth, drop the log
%      from the figure entirely and show only the three annotated views —
%      the inference figures are already stated in the prose.
%
%  OPTIONAL, if a fourth image is wanted:
%      testbed_hardware.jpg — a photograph of the laptop, the USB gigabit
%      adapter, the single Ethernet cable and the Pi 4B. This is the only
%      image that shows the work is hardware-in-the-loop rather than pure
%      simulation. It would go in Section \ref{sec:test-env}.
% =====================================================================

%
%  Intended to follow % =====================================================================
%  Hardware-in-the-Loop Edge Processing
%  A SECTION of the FYP report. Include from the parent chapter with:
%      % =====================================================================
%  Hardware-in-the-Loop Edge Processing
%  A SECTION of the FYP report. Include from the parent chapter with:
%      \input{hitl_edge_processing}
%
%  HEADING LEVELS USED
%      \section     - the HITL section itself
%      \subsection  - its parts
%      \paragraph   - short findings within Results
%
%  REQUIRED PACKAGES (add to the main preamble if not already present)
%      \usepackage{booktabs}
%      \usepackage{graphicx}
%      \usepackage{tikz}
%      \usetikzlibrary{arrows.meta, positioning, fit, backgrounds}
%
%  All labels are prefixed hitl / tab:hitl- / fig:hitl- so they cannot
%  collide with labels elsewhere in the shared report.
% =====================================================================

\section{Hardware-in-the-Loop Edge Processing}
\label{sec:hitl}

The vision workload was moved off the simulation host and onto physical
hardware, so that the cost of onboard detection could be measured on a processor
representative of one that would actually fly. A Raspberry~Pi~4B performs
detection on imagery produced by the simulation, and returns its results across a
wireless channel simulated in ns-3.

% ---------------------------------------------------------------------
\subsection{Extending the testbed with physical hardware}
\label{sec:hitl-testbed}

The existing stack ran entirely on one machine: Gazebo for the environment and
sensors, ArduPilot SITL for the autopilot, ns-3 for the wireless channel, and
Linux network namespaces isolating the ground station from each aircraft. The
ns-3 scenario defined four nodes --- a ground control station and three aircraft
--- of which the second and third were unused placeholders.

A Raspberry~Pi~4B was introduced as the physical edge node and assigned to the
second aircraft slot. Its TAP device, previously created only to satisfy the
scenario's argument list, was brought into service as the attachment point. The
autopilot, physics and channel simulation remained on the host; only the vision
computation moved onto hardware.

The Pi was configured with Ubuntu~22.04 and ROS\,2 Humble to match the host
exactly, giving both machines the same DDS implementation. The host has no
Ethernet port, so a USB~3.0 gigabit adapter was fitted. An initial low-cost
adapter was replaced after \texttt{iperf3} measured \(4.5\)\,Mbps with \(24\%\)
loss, despite \texttt{ethtool} reporting \(100\)\,Mbps full duplex; the
replacement sustained \(834\)\,Mbps with no retransmissions.

% ---------------------------------------------------------------------
\subsection{Constructing the two-link network}
\label{sec:hitl-network}

Two logically separate paths were required between host and Pi. One carries
camera imagery and must remain unimpaired, since on a real aircraft it is a short
cable inside a single airframe rather than a radio. The other carries detection
results and must pass through the simulated channel, as that is the link the
experiment exists to characterise. The same reasoning applies to the
flight-dynamics traffic between SITL and Gazebo, which also bypasses ns-3.

The Pi provides a single Ethernet interface, so both paths were carried over one
cable using IEEE~802.1Q VLAN tagging. On the host, two virtual interfaces were
created on the physical adapter: \texttt{eth-cam}, tagged VLAN~10 and addressed
\(10.0.0.1\); and \texttt{eth-rf}, tagged VLAN~42, left without an address and
made a member of a new bridge, \texttt{br-uav2}. That bridge also received
\texttt{tap-uav2}, the ns-3 interface for the second aircraft node. On the Pi,
matching interfaces \texttt{eth0.10} (\(10.0.0.2\)) and \texttt{eth0.42}
(\(10.42.0.12\)) were created.

Attaching the physical interface to the bridge directly was rejected during
design. A Linux bridge floods frames to all member ports, so camera traffic would
have been delivered into the ns-3 TAP device as well, injecting approximately
\(111\)\,Mbps of imagery into the simulated channel.

The resulting arrangement enforces separation through routing rather than
filtering. Gazebo is addressed only on \(10.0.0.1\) and can reach the Pi only at
\(10.0.0.2\); the ground station namespace is addressed only on \(10.42.0.10\)
and can reach it only at \(10.42.0.12\). Neither can use the other's path.

\begin{figure}[htbp]
\centering
\begin{tikzpicture}[
  font=\small,
  box/.style={draw, rounded corners=2pt, align=center, inner sep=5pt},
  proc/.style={box, fill=black!4},
  net/.style={box, fill=black!10},
  arr/.style={-{Latex[length=2mm]}, thick},
]
\node[proc] (gz)   {Gazebo\\camera + physics};
\node[proc, below=14mm of gz] (gcs) {Ground station\\(namespace)\\\texttt{10.42.0.10}};
\node[net,  right=16mm of gcs] (ns3) {ns-3\\channel};
\node[net,  above=8mm of ns3]  (br)  {\texttt{br-uav2}\\+ TAP};
\node[proc, right=46mm of gz] (pi) {Raspberry Pi 4B\\YOLOv8n detector};

\draw[arr] (gz) -- node[above, align=center, font=\scriptsize]
  {VLAN 10 \quad camera, 2.76\,MB/frame\\unimpaired} (pi);
\draw[arr] (pi.south) |- node[pos=0.75, above, font=\scriptsize]
  {VLAN 42 \quad detections, 118\,B} (br.east);
\draw[arr] (br) -- (ns3);
\draw[arr] (ns3) -- (gcs);

\begin{scope}[on background layer]
  \node[draw, dashed, rounded corners, fit=(gz)(gcs)(ns3)(br),
        inner sep=7pt, label={[font=\scriptsize\itshape]above left:Host PC}] {};
\end{scope}
\end{tikzpicture}
\caption{The two-link arrangement. Camera imagery reaches the edge node over an
unimpaired wired path; detection results traverse the simulated channel. Both
share one physical cable, separated by VLAN tagging.}
\label{fig:hitl-architecture}
\end{figure}

Addressing was made persistent using NetworkManager on the host and netplan on
the Pi, the two distributions using different network managers. Addresses
assigned imperatively were found to be removed within minutes, as each manager
reconciles kernel state against its own configuration and deletes entries it did
not create.

% ---------------------------------------------------------------------
\subsection{Configuring the middleware}
\label{sec:hitl-dds}

A Fast DDS profile was written and applied on both machines. It pins participants
to the intended interfaces through an interface whitelist, since both machines
are attached to WiFi as well as the wired link and DDS otherwise advertises on
every available interface. The profile also disables the shared-memory transport,
which is not isolated by network namespaces and would allow participants to
exchange data without traversing the simulated channel.

Kernel socket buffers were raised on both machines from the \(208\)\,kB default
to \(512\)\,MB. A single \(1280 \times 720\) frame is \(2.76\)\,MB and fragments
into approximately \(1{,}900\) UDP datagrams, of which the default buffer holds
about \(7\%\); reassembly never completed and subscribers received nothing, with
no error reported.

The camera subscription on the Pi was changed from the ROS\,2 default of
\texttt{RELIABLE} with queue depth ten to \texttt{BEST\_EFFORT} with depth one.
Under the default policy the publisher retains each sample until the subscriber
acknowledges it, and because the Pi acknowledges at approximately \(1\)\,Hz, the
camera output of the entire simulation was throttled to \(0.27\)\,Hz.
\texttt{RELIABLE} was retained for detection messages, which must not be lost.

% ---------------------------------------------------------------------
\subsection{Deploying the detector}
\label{sec:hitl-detector}

The detection node runs YOLOv8n restricted to the COCO \emph{person} class,
subscribing to the camera topic and publishing results as JSON. It was deployed
to the Pi and instrumented to record inference duration, measured around the
model invocation alone so that figures remain comparable between back-ends and
against standalone execution.

Two dependency conflicts were resolved during deployment. Installing
\texttt{ultralytics} brings in NumPy~2.x, against which \texttt{cv\_bridge} --- a
compiled extension built for the 1.x C ABI --- fails. Installing an older OpenCV
to compensate reintroduces NumPy~2 as a dependency, and OpenCV~5.0 separately
breaks \texttt{cv\_bridge} by renaming the type constants it uses to build its
conversion table. Versions were pinned to \texttt{numpy==1.26.4} and
\texttt{opencv-python==4.10.0.84}, the only combination satisfying both.

The model was additionally exported to NCNN for ARM-optimised inference. Because
NCNN fixes the input tensor shape at export time, the export geometry was
specified explicitly as \(544 \times 960\) to match the letterboxed frame; a
default export produces a square input in which approximately \(43\%\) of the
computation is spent on padding, at \(19.9\)\,GFLOPS against \(11.3\)\,GFLOPS for
the correct shape. Half-precision export was avoided, the Cortex-A72 lacking fp16
SIMD support.

% ---------------------------------------------------------------------
\subsection{Automation}
\label{sec:hitl-automation}

Bring-up was consolidated into a single script that starts the namespaces, ns-3,
Gazebo, SITL, the micro-ROS agent and the ground station receiver, creates the
VLANs and bridge, and launches the detector on the Pi over SSH. Pre-flight checks
were added for each failure encountered during development: presence of the host
address before Gazebo starts, since Fast DDS enumerates interfaces only at
participant creation and cannot use one assigned later; reachability of the Pi;
adequate socket buffer limits on both machines; and confirmation after startup
that Gazebo bound the intended interface.

% ---------------------------------------------------------------------
\subsection{Results}
\label{sec:hitl-results}

Detections were verified visually, by inspecting annotated frames, before any
quantitative result was accepted. An earlier configuration had reported up to
sixteen detections against five subjects; these were false positives on building
texture, eliminated by raising the confidence threshold and narrowing the camera
field of view.

\paragraph{Channel characterisation.}
Table~\ref{tab:hitl-links} compares the two paths, each measured over thirty ICMP
samples. Short samples are unreliable: the first packet includes address
resolution across the simulated channel, and a two-sample measurement of the same
configuration returned a mean of \(167\)\,ms against \(59\)\,ms for thirty.

\begin{table}[htbp]
\centering
\caption{Measured characteristics of the two links, thirty samples each.}
\label{tab:hitl-links}
\begin{tabular}{lrrrr}
\toprule
Path & Min (ms) & Mean (ms) & Max (ms) & Jitter (ms) \\
\midrule
Sensor, VLAN 10 (wired)   & 0.4  & 1.5  & 2.7   & 0.4  \\
Radio, VLAN 42 (via ns-3) & 11.7 & 77.4 & 242.6 & 54.8 \\
\bottomrule
\end{tabular}
\end{table}

Packet loss was zero on both paths, the aircraft remaining close to the ground
station throughout. The radio path is approximately \(53\times\) slower than the
wired path on average and varies approximately \(140\times\) more.

The distribution within the radio path is itself informative. Its minimum of
\(11.7\)\,ms against a mean of \(77.4\)\,ms is a spread of some \(6.6\times\) on
one link, where a channel applying a fixed delay would produce a narrow one.
The variation indicates that ns-3 is modelling contention, in which packets
queue and back off before transmission, rather than simply adding a constant
offset to each packet.

\paragraph{Sensor transport.}
With the camera at \(20\)\,Hz the wired path carried \(481\)\,Mbps at
\(42{,}544\) packets per second with zero dropped packets, occupying
approximately \(58\%\) of measured link capacity.

\paragraph{Inference latency.}
Table~\ref{tab:hitl-inference} reports latency for each combination of back-end
and input resolution. Changing back-end at constant resolution gives
approximately \(2.5\times\); halving the pixel count at constant back-end gives
approximately \(2.3\times\). Latency therefore scales close to linearly with
pixel count, a pixel ratio of \(0.47\) producing a time ratio of \(0.43\).

\begin{table}[htbp]
\centering
\caption{Inference latency on the Raspberry Pi 4B.}
\label{tab:hitl-inference}
\begin{tabular}{llrr}
\toprule
Back-end & Input & Latency (ms) & Rate (fps) \\
\midrule
PyTorch & $960 \times 544$ & 2540 & 0.39 \\
PyTorch & $640 \times 384$ & 1100 & 0.91 \\
NCNN    & $960 \times 544$ & 1000 & 1.00 \\
NCNN    & $640 \times 384$ & \multicolumn{2}{l}{not measured (est.\ 470)} \\
\bottomrule
\end{tabular}
\end{table}

\paragraph{Detection performance.}
Between two and four of the five subjects were detected per frame, a recall of
\(40\)--\(80\%\). The shortfall is attributable to scene geometry rather than
detector deficiency: the models are separated by approximately one metre, so from
the oblique camera view the nearer figures partially occlude those behind them,
and partially occluded subjects at approximately \(53\) pixels in height fall
below the confidence threshold. Because the number and position of subjects are
known exactly from the world definition, recall is a directly measurable quantity
rather than an estimate.

\paragraph{Transmitted volume.}
Each detection message is approximately \(118\)\,B, against \(2.76\)\,MB for the
corresponding frame, a reduction of approximately \(23{,}000\times\) in data
crossing the radio.

% ---------------------------------------------------------------------
\subsection{Contention between sensor transport and inference}
\label{sec:hitl-contention}

The two workloads on the edge node proved not to be independent. The same model,
on the same input, required \(623\)\,ms executed on a static image but
\(1{,}343\)\,ms within the live pipeline: a difference of \(720\)\,ms, or
approximately half the available inference throughput.

The edge node received frames at \(19.8\)\,Hz while completing inference at
approximately \(0.74\)\,Hz, discarding roughly \(96\%\) of what it received.
Discarded frames are not free: each requires reassembly of some \(2{,}150\) UDP
fragments, deserialisation of \(2.76\)\,MB and a further \(2.76\)\,MB
colour-space conversion, all competing with inference for the same four
processor cores.

Reducing the camera rate at source from \(20\)\,Hz to \(5\)\,Hz --- still
approximately seven times the rate the detector can consume --- reduced inference
latency to approximately \(1{,}000\)\,ms, recovering about a quarter of the
deficit, and reduced link occupancy from \(481\)\,Mbps to \(111\)\,Mbps. The
application-level frame-rate control provided by the relay node cannot achieve
this, as it discards frames after reception, by which point the transport and
deserialisation cost has already been incurred.

This interaction is not observable in pure simulation, where sensor data is
delivered through shared memory at negligible cost, nor in an isolated benchmark
of the detector, since the interference originates outside it. It appears only
when a physical processor performs both reception and inference concurrently.

% ---------------------------------------------------------------------
\subsection{Limitations}
\label{sec:hitl-limitations}

Clock synchronisation between the two machines has not yet been established.
End-to-end latency is computed from timestamps taken on different hosts, and
without synchronisation those values are not meaningful while remaining plausible
in magnitude; synchronisation is required before end-to-end latency is reported.

The edge configuration is characterised; the corresponding ground configuration,
in which compressed imagery crosses the simulated channel and inference is
performed at the ground station, remains to be measured.

Measurements were taken with the aircraft close to the ground station, giving
zero packet loss throughout. Behaviour as range increases and buildings obstruct
the path has not been characterised. The testbed supports a single physical edge
node; multi-node operation has not been evaluated.

%
%  HEADING LEVELS USED
%      \section     - the HITL section itself
%      \subsection  - its parts
%      \paragraph   - short findings within Results
%
%  REQUIRED PACKAGES (add to the main preamble if not already present)
%      \usepackage{booktabs}
%      \usepackage{graphicx}
%      \usepackage{tikz}
%      \usetikzlibrary{arrows.meta, positioning, fit, backgrounds}
%
%  All labels are prefixed hitl / tab:hitl- / fig:hitl- so they cannot
%  collide with labels elsewhere in the shared report.
% =====================================================================

\section{Hardware-in-the-Loop Edge Processing}
\label{sec:hitl}

The vision workload was moved off the simulation host and onto physical
hardware, so that the cost of onboard detection could be measured on a processor
representative of one that would actually fly. A Raspberry~Pi~4B performs
detection on imagery produced by the simulation, and returns its results across a
wireless channel simulated in ns-3.

% ---------------------------------------------------------------------
\subsection{Extending the testbed with physical hardware}
\label{sec:hitl-testbed}

The existing stack ran entirely on one machine: Gazebo for the environment and
sensors, ArduPilot SITL for the autopilot, ns-3 for the wireless channel, and
Linux network namespaces isolating the ground station from each aircraft. The
ns-3 scenario defined four nodes --- a ground control station and three aircraft
--- of which the second and third were unused placeholders.

A Raspberry~Pi~4B was introduced as the physical edge node and assigned to the
second aircraft slot. Its TAP device, previously created only to satisfy the
scenario's argument list, was brought into service as the attachment point. The
autopilot, physics and channel simulation remained on the host; only the vision
computation moved onto hardware.

The Pi was configured with Ubuntu~22.04 and ROS\,2 Humble to match the host
exactly, giving both machines the same DDS implementation. The host has no
Ethernet port, so a USB~3.0 gigabit adapter was fitted. An initial low-cost
adapter was replaced after \texttt{iperf3} measured \(4.5\)\,Mbps with \(24\%\)
loss, despite \texttt{ethtool} reporting \(100\)\,Mbps full duplex; the
replacement sustained \(834\)\,Mbps with no retransmissions.

% ---------------------------------------------------------------------
\subsection{Constructing the two-link network}
\label{sec:hitl-network}

Two logically separate paths were required between host and Pi. One carries
camera imagery and must remain unimpaired, since on a real aircraft it is a short
cable inside a single airframe rather than a radio. The other carries detection
results and must pass through the simulated channel, as that is the link the
experiment exists to characterise. The same reasoning applies to the
flight-dynamics traffic between SITL and Gazebo, which also bypasses ns-3.

The Pi provides a single Ethernet interface, so both paths were carried over one
cable using IEEE~802.1Q VLAN tagging. On the host, two virtual interfaces were
created on the physical adapter: \texttt{eth-cam}, tagged VLAN~10 and addressed
\(10.0.0.1\); and \texttt{eth-rf}, tagged VLAN~42, left without an address and
made a member of a new bridge, \texttt{br-uav2}. That bridge also received
\texttt{tap-uav2}, the ns-3 interface for the second aircraft node. On the Pi,
matching interfaces \texttt{eth0.10} (\(10.0.0.2\)) and \texttt{eth0.42}
(\(10.42.0.12\)) were created.

Attaching the physical interface to the bridge directly was rejected during
design. A Linux bridge floods frames to all member ports, so camera traffic would
have been delivered into the ns-3 TAP device as well, injecting approximately
\(111\)\,Mbps of imagery into the simulated channel.

The resulting arrangement enforces separation through routing rather than
filtering. Gazebo is addressed only on \(10.0.0.1\) and can reach the Pi only at
\(10.0.0.2\); the ground station namespace is addressed only on \(10.42.0.10\)
and can reach it only at \(10.42.0.12\). Neither can use the other's path.

\begin{figure}[htbp]
\centering
\begin{tikzpicture}[
  font=\small,
  box/.style={draw, rounded corners=2pt, align=center, inner sep=5pt},
  proc/.style={box, fill=black!4},
  net/.style={box, fill=black!10},
  arr/.style={-{Latex[length=2mm]}, thick},
]
\node[proc] (gz)   {Gazebo\\camera + physics};
\node[proc, below=14mm of gz] (gcs) {Ground station\\(namespace)\\\texttt{10.42.0.10}};
\node[net,  right=16mm of gcs] (ns3) {ns-3\\channel};
\node[net,  above=8mm of ns3]  (br)  {\texttt{br-uav2}\\+ TAP};
\node[proc, right=46mm of gz] (pi) {Raspberry Pi 4B\\YOLOv8n detector};

\draw[arr] (gz) -- node[above, align=center, font=\scriptsize]
  {VLAN 10 \quad camera, 2.76\,MB/frame\\unimpaired} (pi);
\draw[arr] (pi.south) |- node[pos=0.75, above, font=\scriptsize]
  {VLAN 42 \quad detections, 118\,B} (br.east);
\draw[arr] (br) -- (ns3);
\draw[arr] (ns3) -- (gcs);

\begin{scope}[on background layer]
  \node[draw, dashed, rounded corners, fit=(gz)(gcs)(ns3)(br),
        inner sep=7pt, label={[font=\scriptsize\itshape]above left:Host PC}] {};
\end{scope}
\end{tikzpicture}
\caption{The two-link arrangement. Camera imagery reaches the edge node over an
unimpaired wired path; detection results traverse the simulated channel. Both
share one physical cable, separated by VLAN tagging.}
\label{fig:hitl-architecture}
\end{figure}

Addressing was made persistent using NetworkManager on the host and netplan on
the Pi, the two distributions using different network managers. Addresses
assigned imperatively were found to be removed within minutes, as each manager
reconciles kernel state against its own configuration and deletes entries it did
not create.

% ---------------------------------------------------------------------
\subsection{Configuring the middleware}
\label{sec:hitl-dds}

A Fast DDS profile was written and applied on both machines. It pins participants
to the intended interfaces through an interface whitelist, since both machines
are attached to WiFi as well as the wired link and DDS otherwise advertises on
every available interface. The profile also disables the shared-memory transport,
which is not isolated by network namespaces and would allow participants to
exchange data without traversing the simulated channel.

Kernel socket buffers were raised on both machines from the \(208\)\,kB default
to \(512\)\,MB. A single \(1280 \times 720\) frame is \(2.76\)\,MB and fragments
into approximately \(1{,}900\) UDP datagrams, of which the default buffer holds
about \(7\%\); reassembly never completed and subscribers received nothing, with
no error reported.

The camera subscription on the Pi was changed from the ROS\,2 default of
\texttt{RELIABLE} with queue depth ten to \texttt{BEST\_EFFORT} with depth one.
Under the default policy the publisher retains each sample until the subscriber
acknowledges it, and because the Pi acknowledges at approximately \(1\)\,Hz, the
camera output of the entire simulation was throttled to \(0.27\)\,Hz.
\texttt{RELIABLE} was retained for detection messages, which must not be lost.

% ---------------------------------------------------------------------
\subsection{Deploying the detector}
\label{sec:hitl-detector}

The detection node runs YOLOv8n restricted to the COCO \emph{person} class,
subscribing to the camera topic and publishing results as JSON. It was deployed
to the Pi and instrumented to record inference duration, measured around the
model invocation alone so that figures remain comparable between back-ends and
against standalone execution.

Two dependency conflicts were resolved during deployment. Installing
\texttt{ultralytics} brings in NumPy~2.x, against which \texttt{cv\_bridge} --- a
compiled extension built for the 1.x C ABI --- fails. Installing an older OpenCV
to compensate reintroduces NumPy~2 as a dependency, and OpenCV~5.0 separately
breaks \texttt{cv\_bridge} by renaming the type constants it uses to build its
conversion table. Versions were pinned to \texttt{numpy==1.26.4} and
\texttt{opencv-python==4.10.0.84}, the only combination satisfying both.

The model was additionally exported to NCNN for ARM-optimised inference. Because
NCNN fixes the input tensor shape at export time, the export geometry was
specified explicitly as \(544 \times 960\) to match the letterboxed frame; a
default export produces a square input in which approximately \(43\%\) of the
computation is spent on padding, at \(19.9\)\,GFLOPS against \(11.3\)\,GFLOPS for
the correct shape. Half-precision export was avoided, the Cortex-A72 lacking fp16
SIMD support.

% ---------------------------------------------------------------------
\subsection{Automation}
\label{sec:hitl-automation}

Bring-up was consolidated into a single script that starts the namespaces, ns-3,
Gazebo, SITL, the micro-ROS agent and the ground station receiver, creates the
VLANs and bridge, and launches the detector on the Pi over SSH. Pre-flight checks
were added for each failure encountered during development: presence of the host
address before Gazebo starts, since Fast DDS enumerates interfaces only at
participant creation and cannot use one assigned later; reachability of the Pi;
adequate socket buffer limits on both machines; and confirmation after startup
that Gazebo bound the intended interface.

% ---------------------------------------------------------------------
\subsection{Results}
\label{sec:hitl-results}

Detections were verified visually, by inspecting annotated frames, before any
quantitative result was accepted. An earlier configuration had reported up to
sixteen detections against five subjects; these were false positives on building
texture, eliminated by raising the confidence threshold and narrowing the camera
field of view.

\paragraph{Channel characterisation.}
Table~\ref{tab:hitl-links} compares the two paths, each measured over thirty ICMP
samples. Short samples are unreliable: the first packet includes address
resolution across the simulated channel, and a two-sample measurement of the same
configuration returned a mean of \(167\)\,ms against \(59\)\,ms for thirty.

\begin{table}[htbp]
\centering
\caption{Measured characteristics of the two links, thirty samples each.}
\label{tab:hitl-links}
\begin{tabular}{lrrrr}
\toprule
Path & Min (ms) & Mean (ms) & Max (ms) & Jitter (ms) \\
\midrule
Sensor, VLAN 10 (wired)   & 0.4  & 1.5  & 2.7   & 0.4  \\
Radio, VLAN 42 (via ns-3) & 11.7 & 77.4 & 242.6 & 54.8 \\
\bottomrule
\end{tabular}
\end{table}

Packet loss was zero on both paths, the aircraft remaining close to the ground
station throughout. The radio path is approximately \(53\times\) slower than the
wired path on average and varies approximately \(140\times\) more.

The distribution within the radio path is itself informative. Its minimum of
\(11.7\)\,ms against a mean of \(77.4\)\,ms is a spread of some \(6.6\times\) on
one link, where a channel applying a fixed delay would produce a narrow one.
The variation indicates that ns-3 is modelling contention, in which packets
queue and back off before transmission, rather than simply adding a constant
offset to each packet.

\paragraph{Sensor transport.}
With the camera at \(20\)\,Hz the wired path carried \(481\)\,Mbps at
\(42{,}544\) packets per second with zero dropped packets, occupying
approximately \(58\%\) of measured link capacity.

\paragraph{Inference latency.}
Table~\ref{tab:hitl-inference} reports latency for each combination of back-end
and input resolution. Changing back-end at constant resolution gives
approximately \(2.5\times\); halving the pixel count at constant back-end gives
approximately \(2.3\times\). Latency therefore scales close to linearly with
pixel count, a pixel ratio of \(0.47\) producing a time ratio of \(0.43\).

\begin{table}[htbp]
\centering
\caption{Inference latency on the Raspberry Pi 4B.}
\label{tab:hitl-inference}
\begin{tabular}{llrr}
\toprule
Back-end & Input & Latency (ms) & Rate (fps) \\
\midrule
PyTorch & $960 \times 544$ & 2540 & 0.39 \\
PyTorch & $640 \times 384$ & 1100 & 0.91 \\
NCNN    & $960 \times 544$ & 1000 & 1.00 \\
NCNN    & $640 \times 384$ & \multicolumn{2}{l}{not measured (est.\ 470)} \\
\bottomrule
\end{tabular}
\end{table}

\paragraph{Detection performance.}
Between two and four of the five subjects were detected per frame, a recall of
\(40\)--\(80\%\). The shortfall is attributable to scene geometry rather than
detector deficiency: the models are separated by approximately one metre, so from
the oblique camera view the nearer figures partially occlude those behind them,
and partially occluded subjects at approximately \(53\) pixels in height fall
below the confidence threshold. Because the number and position of subjects are
known exactly from the world definition, recall is a directly measurable quantity
rather than an estimate.

\paragraph{Transmitted volume.}
Each detection message is approximately \(118\)\,B, against \(2.76\)\,MB for the
corresponding frame, a reduction of approximately \(23{,}000\times\) in data
crossing the radio.

% ---------------------------------------------------------------------
\subsection{Contention between sensor transport and inference}
\label{sec:hitl-contention}

The two workloads on the edge node proved not to be independent. The same model,
on the same input, required \(623\)\,ms executed on a static image but
\(1{,}343\)\,ms within the live pipeline: a difference of \(720\)\,ms, or
approximately half the available inference throughput.

The edge node received frames at \(19.8\)\,Hz while completing inference at
approximately \(0.74\)\,Hz, discarding roughly \(96\%\) of what it received.
Discarded frames are not free: each requires reassembly of some \(2{,}150\) UDP
fragments, deserialisation of \(2.76\)\,MB and a further \(2.76\)\,MB
colour-space conversion, all competing with inference for the same four
processor cores.

Reducing the camera rate at source from \(20\)\,Hz to \(5\)\,Hz --- still
approximately seven times the rate the detector can consume --- reduced inference
latency to approximately \(1{,}000\)\,ms, recovering about a quarter of the
deficit, and reduced link occupancy from \(481\)\,Mbps to \(111\)\,Mbps. The
application-level frame-rate control provided by the relay node cannot achieve
this, as it discards frames after reception, by which point the transport and
deserialisation cost has already been incurred.

This interaction is not observable in pure simulation, where sensor data is
delivered through shared memory at negligible cost, nor in an isolated benchmark
of the detector, since the interference originates outside it. It appears only
when a physical processor performs both reception and inference concurrently.

% ---------------------------------------------------------------------
\subsection{Limitations}
\label{sec:hitl-limitations}

Clock synchronisation between the two machines has not yet been established.
End-to-end latency is computed from timestamps taken on different hosts, and
without synchronisation those values are not meaningful while remaining plausible
in magnitude; synchronisation is required before end-to-end latency is reported.

The edge configuration is characterised; the corresponding ground configuration,
in which compressed imagery crosses the simulated channel and inference is
performed at the ground station, remains to be measured.

Measurements were taken with the aircraft close to the ground station, giving
zero packet loss throughout. Behaviour as range increases and buildings obstruct
the path has not been characterised. The testbed supports a single physical edge
node; multi-node operation has not been evaluated.
; cross-references
%  \ref{sec:hitl}. A longer version with the full 26-case test matrix is
%  kept in testing_validation_long.tex if an appendix is wanted.
%
%  REQUIRED PACKAGES
%      \usepackage{booktabs}
%      \usepackage{array}
%      \usepackage{graphicx}
%      \usepackage{subcaption}
%      \graphicspath{{figures/}}
%
%  IMAGES EXPECTED in docs/report/figures/  (see checklist at end)
%      detection_a.png  detection_b.png  detection_c.png
% =====================================================================

\section{Testing, Experimentation and Validation}
\label{sec:test}

Testing of the hardware-in-the-loop testbed described in
Section~\ref{sec:hitl} was organised around a difficulty specific to this
kind of system: a fault at a low layer does not raise an error at the layer
above it, but produces plausible and wrong behaviour there. A subscriber
starved of buffer space receives nothing and reports nothing; two nodes with
mismatched topic names run indefinitely without complaint. Testing was
therefore performed bottom-up --- physical link, path separation, middleware
transport, detector, end-to-end --- with each layer required to pass before
the next was attempted.

Three results are reported here. The first establishes that the two-link
architecture does what it claims. The second establishes that the detector
finds real people rather than plausible-looking texture. The third concerns
the measurement chain itself, and is included because it determines whether
the other two can be believed.

% ---------------------------------------------------------------------
\subsection{Test environment}
\label{sec:test-env}

Both machines run Ubuntu~22.04 with ROS\,2 Humble and Fast~DDS~2.6, so the
middleware is identical at each end of the link and cannot itself be a
source of difference. The simulation host runs Gazebo Classic in lockstep
with ArduPilot SITL, and ns-3 with TapBridge for the wireless channel. The
edge node is a Raspberry~Pi~4B running YOLOv8n restricted to the COCO
\emph{person} class. The camera produces $1280 \times 720$ frames of
$2.76$\,MB at $5$\,Hz. The test scene contains exactly five human models at
coordinates fixed in the world definition, which makes recall a counted
quantity rather than an estimated one.

The Pi's single Ethernet port is divided by 802.1Q tagging into
\texttt{eth0.10} ($10.0.0.2$, camera, unimpaired) and \texttt{eth0.42}
($10.42.0.12$, detections, through ns-3), created by
\texttt{scripts/netns/pi\_hitl\_link.sh}. The matching host interfaces
\texttt{eth-cam} and \texttt{eth-rf} are created by step~1c of
\texttt{scripts/netns/launch\_single\_uav\_netns.sh}, the latter bridged
into \texttt{tap-uav2}.

% ---------------------------------------------------------------------
\subsection{Result 1: the two links behave differently}
\label{sec:test-links}

The central claim of the architecture is that detection results cross a
simulated wireless channel while camera imagery does not. Two paths that
both work are not evidence of this; they must also be shown to differ in the
intended way, and imagery must be shown to be absent from the radio path.

Table~\ref{tab:test-links} reports sixty ICMP samples on the sensor path and
thirty on the radio path. Large samples are necessary: the first packet on the
radio path includes address resolution across the simulated channel, and a
two-sample measurement of the same configuration returned a mean of
$167$\,ms against $74.5$\,ms for thirty. The short sample is biased, not
merely noisy.

\begin{table}[htbp]
\centering
\caption{Measured characteristics of the two paths. Sensor path over sixty
samples, radio path over thirty; zero packet loss on both.}
\label{tab:test-links}
\begin{tabular}{@{}lrrrr@{}}
\toprule
Path & Min (ms) & Mean (ms) & Max (ms) & Jitter (ms) \\
\midrule
Sensor, VLAN 10 (wired)   & 1.0  & 1.4  & 1.8   & 0.12 \\
Radio, VLAN 42 (via ns-3) & 10.1 & 74.5 & 158.0 & 38.4 \\
\midrule
Ratio                     & --   & $53\times$ & -- & $320\times$ \\
\bottomrule
\end{tabular}
\end{table}

The radio path is approximately \(53\times\) slower on average, and its
variation is larger by more than two orders of magnitude.

Its internal spread is itself informative: a minimum of \(10.1\)\,ms against
a mean of \(74.5\)\,ms is a range of some \(7\times\) on one link, where a
channel applying a fixed delay would produce a narrow distribution. The
variation indicates that ns-3 is modelling contention, in which packets queue
and back off, rather than adding a constant offset.

The radio measurement was repeated. An independent set of thirty samples taken
several hours earlier, with the host reconfigured in between, gave a mean of
\(77.4\)\,ms against the \(74.5\)\,ms reported here, agreeing within
\(4\%\).

The sensor figures were taken \emph{while} the camera was streaming at
\(110\)\,Mbps. An idle-cable measurement over the same number of samples
gives the same mean, so the two logical links sharing one physical cable do
not delay one another. Queue management was verified rather than assumed:
both machines run \texttt{fq\_codel} with a \(5\)\,ms target, and byte
queue limits on the edge node had reduced its driver ring to approximately one
packet.

Absence of imagery on the radio path was established negatively, by
observation rather than by inspecting a rule set, since separation here is
enforced by addressing and bridge membership rather than by filtering. No
$10.0.0.0/24$ traffic was observed on \texttt{tap-uav2}, the point at which
traffic enters ns-3, and the byte counters on the two interfaces differ by
approximately four orders of magnitude --- consistent with $2.76$\,MB per
frame on one path against $118$\,B per detection message on the other, a
reduction of approximately \(23{,}000\times\).

A structural argument supports the measurements. The ground station receiver
executes inside the \texttt{gcsns} network namespace, whose only interface
is one end of a veth pair attached to \texttt{tap-gcs}. That namespace has
no interface on $10.0.0.0/24$ and no route to it. A message arriving at the
receiver cannot have taken any other path, because no other path exists from
inside the namespace. Separation is therefore a property of the topology
rather than of a setting that could be silently overridden.

% ---------------------------------------------------------------------
\subsection{Result 2: detections correspond to real subjects}
\label{sec:test-detection}

Visual confirmation was made a precondition of accepting any quantitative
result. A detector reporting a plausible count while placing its boxes on
building texture would satisfy every numerical check in this section, and an
earlier configuration did exactly that: sixteen confident detections against
five subjects, none correct. Only inspection of annotated frames
distinguishes the two cases. Those false positives were eliminated by
raising the confidence threshold and narrowing the camera field of view.

\begin{figure}[htbp]
\centering
\begin{subfigure}[b]{0.32\textwidth}
  \includegraphics[width=\textwidth]{detection_a.png}
  \caption{Two subjects, confidence $0.59$ and $0.42$.}
  \label{fig:test-det-a}
\end{subfigure}
\hfill
\begin{subfigure}[b]{0.32\textwidth}
  \includegraphics[width=\textwidth]{detection_b.png}
  \caption{One subject, confidence $0.48$.}
  \label{fig:test-det-b}
\end{subfigure}
\hfill
\begin{subfigure}[b]{0.32\textwidth}
  \includegraphics[width=\textwidth]{detection_c.png}
  \caption{Three subjects, confidence $0.59$, $0.47$ and $0.46$.}
  \label{fig:test-det-c}
\end{subfigure}
\caption{Detector output during flight, with the edge node's log above each
view. Bounding boxes coincide with the human models rather than with ground
texture. The log records every frame, including those with no detection, so
that a stalled node is distinguishable from an empty scene, and reports
inference duration per frame.}
\label{fig:test-detection}
\end{figure}

Figure~\ref{fig:test-detection} shows three moments from a patrol flight.
Between one and four of the five subjects were detected per frame, at
confidences of \(0.42\)--\(0.59\). The shortfall against the known count of
five is attributable to scene geometry rather than to detector deficiency:
the models stand approximately one metre apart, so from the oblique camera
view the nearer figures partially occlude those behind them, and a partially
occluded subject at approximately \(53\) pixels in height falls below the
confidence threshold. Because the number of subjects is known exactly, any
count exceeding five would be a false positive without the need to
adjudicate individual boxes; none was observed in this configuration.

Inference duration was measured around the model invocation alone, excluding
message reception, format conversion and publication, so that the figure
remains comparable between back-ends and against standalone execution. Over
\(480\) consecutive frames of one patrol flight the mean was
\(1{,}093\)\,ms, with a minimum of \(1{,}025\)\,ms. The distribution is tight
apart from a single outlier at \(6{,}949\)\,ms, one frame in \(480\),
attributable to contention for the processor from another task rather than to
the model itself.

The Pi therefore completes rather less than one detection per second against a
camera delivering five frames per second, and discards approximately \(80\%\)
of what it receives --- which is why the camera subscription uses
\texttt{BEST\_EFFORT} with queue depth one, keeping the newest frame rather
than accumulating a backlog of stale ones.

% ---------------------------------------------------------------------
\subsection{Result 3: validating the measurement chain}
\label{sec:test-chain}

The two results above depend on the camera actually reaching the edge node
at the rate the configuration specifies. That assumption failed, and the
manner of its failure is reported here because it determines what the other
measurements are worth.

The symptom was that the edge node received imagery at \(0.19\)\,Hz where
the camera was configured for \(5\)\,Hz, a shortfall of approximately
\(26\times\). The visible consequences appeared three layers above the
cause: the detector reported no humans in \(93\) of \(94\) frames, and the
patrol mission timed out at every waypoint. Both invite explanation in terms
of detector settings or flight control, and neither has anything to do with
either --- the aircraft flew correctly and the detector functioned
correctly, but the scene was being observed through roughly one frame every
five seconds.

\begin{table}[htbp]
\centering
\caption{Eliminating candidate causes of the frame-rate shortfall. Each row
is a measurement that excludes one layer.}
\label{tab:test-elimination}
\begin{tabular}{@{}p{38mm}p{38mm}p{40mm}@{}}
\toprule
Hypothesis & Measurement & Outcome \\
\midrule
Simulation running slowly
  & Real-time factor from \texttt{/clock}
  & \(0.998\); simulated time tracks wall clock \\[3pt]
Gazebo rendering slowly
  & Host-side subscriber on the camera topic
  & \(5.00\)\,Hz, gaps \(0.19\)--\(0.21\)\,s; source correct \\[3pt]
Link negotiated at $100$\,Mbps
  & \texttt{ethtool} on the edge node
  & $1000$\,Mb/s full duplex \\[3pt]
Fragments lost in flight
  & Interface counters, both machines
  & Zero errors, zero drops \\[3pt]
Receive buffers too small
  & \texttt{net.core.rmem\_max}
  & $512$\,MB, already raised \\[3pt]
\textbf{Send buffers too small}
  & \texttt{net.core.wmem\_max}
  & \(\mathbf{208}\)\,\textbf{kB}, left at the default \\
\bottomrule
\end{tabular}
\end{table}

The DDS profile requests a \(16\)\,MB send buffer. The kernel caps any such
request at \texttt{net.core.wmem\_max}, and performs the reduction silently.
A single frame is therefore approximately \(13\times\) larger than the
buffer available to transmit it: the publisher filled that buffer
immediately and discarded the remaining fragments at the socket layer,
before they reached the network interface. Comparing production with
transmission confirmed this --- the host generated \(414\)\,MB of imagery
over thirty seconds and placed \(10\)\,MB of it on the wire, \(2\)\,Mbps
against the \(110\)\,Mbps the configuration implies. Raising the parameter
on both machines and restarting the stack, which is required because
Fast~DDS applies socket options only at participant creation, restored
delivery to \(5.000\)\,Hz with an inter-arrival standard deviation of
\(4.5\)\,ms: an exact match to the source.

Two properties of this fault are worth recording. First, no counter
registers it. Loss at the socket layer is invisible to interface
statistics, and \texttt{tx\_dropped} on the sending interface and
\texttt{rx\_errors} on the receiving one both remained at zero throughout,
because the discarded fragments were never transmitted. A diagnosis founded
on interface counters would have concluded that the network was healthy,
which it was. Second, the two directions are governed by separate kernel
parameters. The receive path had been corrected during earlier work and the
transmit path had not, because the failure that prompted the original
correction pointed only at the receiving side. The pre-flight checks in
\texttt{scripts/netns/run\_hitl.sh} now verify both parameters on both
machines before the stack starts.

% ---------------------------------------------------------------------
\subsection{Limitations}
\label{sec:test-limitations}

Clock synchronisation between the two machines has not been established.
End-to-end latency would be computed from timestamps taken on different
hosts, and without synchronisation such values are not meaningful while
remaining plausible in magnitude. End-to-end latency is therefore not
reported, rather than reported with a caveat.

The edge configuration is characterised; the corresponding ground
configuration, in which compressed imagery crosses the simulated channel and
inference is performed at the ground station, remains to be measured. All
measurements were taken with the aircraft close to the ground station,
giving zero packet loss on both paths, so the reliability policy applied to
detection messages is untested under the loss conditions it exists to
handle. The testbed supports a single physical edge node; multi-node
operation has not been evaluated.

% =====================================================================
%  IMAGE CHECKLIST — save into docs/report/figures/
%
%  detection_a.png / detection_b.png / detection_c.png
%      The three screenshots already captured: the detector log above and
%      the rqt_image_view annotated frame below, showing 2, 1 and 3
%      detections respectively.
%
%      CROP EACH ONE before including. At three-to-a-row the full desktop
%      screenshot will be illegible. Keep only:
%        - the last ~8 log lines (enough to show the count and the
%          inference time), and
%        - the image pane itself, without the rqt toolbar and topic
%          selector.
%      Crop all three to the same aspect ratio so the row lines up.
%
%      If the log text is still too small at 0.32\textwidth, drop the log
%      from the figure entirely and show only the three annotated views —
%      the inference figures are already stated in the prose.
%
%  OPTIONAL, if a fourth image is wanted:
%      testbed_hardware.jpg — a photograph of the laptop, the USB gigabit
%      adapter, the single Ethernet cable and the Pi 4B. This is the only
%      image that shows the work is hardware-in-the-loop rather than pure
%      simulation. It would go in Section \ref{sec:test-env}.
% =====================================================================

%
%  Intended to follow % =====================================================================
%  Hardware-in-the-Loop Edge Processing
%  A SECTION of the FYP report. Include from the parent chapter with:
%      % =====================================================================
%  Hardware-in-the-Loop Edge Processing
%  A SECTION of the FYP report. Include from the parent chapter with:
%      % =====================================================================
%  Hardware-in-the-Loop Edge Processing
%  A SECTION of the FYP report. Include from the parent chapter with:
%      \input{hitl_edge_processing}
%
%  HEADING LEVELS USED
%      \section     - the HITL section itself
%      \subsection  - its parts
%      \paragraph   - short findings within Results
%
%  REQUIRED PACKAGES (add to the main preamble if not already present)
%      \usepackage{booktabs}
%      \usepackage{graphicx}
%      \usepackage{tikz}
%      \usetikzlibrary{arrows.meta, positioning, fit, backgrounds}
%
%  All labels are prefixed hitl / tab:hitl- / fig:hitl- so they cannot
%  collide with labels elsewhere in the shared report.
% =====================================================================

\section{Hardware-in-the-Loop Edge Processing}
\label{sec:hitl}

The vision workload was moved off the simulation host and onto physical
hardware, so that the cost of onboard detection could be measured on a processor
representative of one that would actually fly. A Raspberry~Pi~4B performs
detection on imagery produced by the simulation, and returns its results across a
wireless channel simulated in ns-3.

% ---------------------------------------------------------------------
\subsection{Extending the testbed with physical hardware}
\label{sec:hitl-testbed}

The existing stack ran entirely on one machine: Gazebo for the environment and
sensors, ArduPilot SITL for the autopilot, ns-3 for the wireless channel, and
Linux network namespaces isolating the ground station from each aircraft. The
ns-3 scenario defined four nodes --- a ground control station and three aircraft
--- of which the second and third were unused placeholders.

A Raspberry~Pi~4B was introduced as the physical edge node and assigned to the
second aircraft slot. Its TAP device, previously created only to satisfy the
scenario's argument list, was brought into service as the attachment point. The
autopilot, physics and channel simulation remained on the host; only the vision
computation moved onto hardware.

The Pi was configured with Ubuntu~22.04 and ROS\,2 Humble to match the host
exactly, giving both machines the same DDS implementation. The host has no
Ethernet port, so a USB~3.0 gigabit adapter was fitted. An initial low-cost
adapter was replaced after \texttt{iperf3} measured \(4.5\)\,Mbps with \(24\%\)
loss, despite \texttt{ethtool} reporting \(100\)\,Mbps full duplex; the
replacement sustained \(834\)\,Mbps with no retransmissions.

% ---------------------------------------------------------------------
\subsection{Constructing the two-link network}
\label{sec:hitl-network}

Two logically separate paths were required between host and Pi. One carries
camera imagery and must remain unimpaired, since on a real aircraft it is a short
cable inside a single airframe rather than a radio. The other carries detection
results and must pass through the simulated channel, as that is the link the
experiment exists to characterise. The same reasoning applies to the
flight-dynamics traffic between SITL and Gazebo, which also bypasses ns-3.

The Pi provides a single Ethernet interface, so both paths were carried over one
cable using IEEE~802.1Q VLAN tagging. On the host, two virtual interfaces were
created on the physical adapter: \texttt{eth-cam}, tagged VLAN~10 and addressed
\(10.0.0.1\); and \texttt{eth-rf}, tagged VLAN~42, left without an address and
made a member of a new bridge, \texttt{br-uav2}. That bridge also received
\texttt{tap-uav2}, the ns-3 interface for the second aircraft node. On the Pi,
matching interfaces \texttt{eth0.10} (\(10.0.0.2\)) and \texttt{eth0.42}
(\(10.42.0.12\)) were created.

Attaching the physical interface to the bridge directly was rejected during
design. A Linux bridge floods frames to all member ports, so camera traffic would
have been delivered into the ns-3 TAP device as well, injecting approximately
\(111\)\,Mbps of imagery into the simulated channel.

The resulting arrangement enforces separation through routing rather than
filtering. Gazebo is addressed only on \(10.0.0.1\) and can reach the Pi only at
\(10.0.0.2\); the ground station namespace is addressed only on \(10.42.0.10\)
and can reach it only at \(10.42.0.12\). Neither can use the other's path.

\begin{figure}[htbp]
\centering
\begin{tikzpicture}[
  font=\small,
  box/.style={draw, rounded corners=2pt, align=center, inner sep=5pt},
  proc/.style={box, fill=black!4},
  net/.style={box, fill=black!10},
  arr/.style={-{Latex[length=2mm]}, thick},
]
\node[proc] (gz)   {Gazebo\\camera + physics};
\node[proc, below=14mm of gz] (gcs) {Ground station\\(namespace)\\\texttt{10.42.0.10}};
\node[net,  right=16mm of gcs] (ns3) {ns-3\\channel};
\node[net,  above=8mm of ns3]  (br)  {\texttt{br-uav2}\\+ TAP};
\node[proc, right=46mm of gz] (pi) {Raspberry Pi 4B\\YOLOv8n detector};

\draw[arr] (gz) -- node[above, align=center, font=\scriptsize]
  {VLAN 10 \quad camera, 2.76\,MB/frame\\unimpaired} (pi);
\draw[arr] (pi.south) |- node[pos=0.75, above, font=\scriptsize]
  {VLAN 42 \quad detections, 118\,B} (br.east);
\draw[arr] (br) -- (ns3);
\draw[arr] (ns3) -- (gcs);

\begin{scope}[on background layer]
  \node[draw, dashed, rounded corners, fit=(gz)(gcs)(ns3)(br),
        inner sep=7pt, label={[font=\scriptsize\itshape]above left:Host PC}] {};
\end{scope}
\end{tikzpicture}
\caption{The two-link arrangement. Camera imagery reaches the edge node over an
unimpaired wired path; detection results traverse the simulated channel. Both
share one physical cable, separated by VLAN tagging.}
\label{fig:hitl-architecture}
\end{figure}

Addressing was made persistent using NetworkManager on the host and netplan on
the Pi, the two distributions using different network managers. Addresses
assigned imperatively were found to be removed within minutes, as each manager
reconciles kernel state against its own configuration and deletes entries it did
not create.

% ---------------------------------------------------------------------
\subsection{Configuring the middleware}
\label{sec:hitl-dds}

A Fast DDS profile was written and applied on both machines. It pins participants
to the intended interfaces through an interface whitelist, since both machines
are attached to WiFi as well as the wired link and DDS otherwise advertises on
every available interface. The profile also disables the shared-memory transport,
which is not isolated by network namespaces and would allow participants to
exchange data without traversing the simulated channel.

Kernel socket buffers were raised on both machines from the \(208\)\,kB default
to \(512\)\,MB. A single \(1280 \times 720\) frame is \(2.76\)\,MB and fragments
into approximately \(1{,}900\) UDP datagrams, of which the default buffer holds
about \(7\%\); reassembly never completed and subscribers received nothing, with
no error reported.

The camera subscription on the Pi was changed from the ROS\,2 default of
\texttt{RELIABLE} with queue depth ten to \texttt{BEST\_EFFORT} with depth one.
Under the default policy the publisher retains each sample until the subscriber
acknowledges it, and because the Pi acknowledges at approximately \(1\)\,Hz, the
camera output of the entire simulation was throttled to \(0.27\)\,Hz.
\texttt{RELIABLE} was retained for detection messages, which must not be lost.

% ---------------------------------------------------------------------
\subsection{Deploying the detector}
\label{sec:hitl-detector}

The detection node runs YOLOv8n restricted to the COCO \emph{person} class,
subscribing to the camera topic and publishing results as JSON. It was deployed
to the Pi and instrumented to record inference duration, measured around the
model invocation alone so that figures remain comparable between back-ends and
against standalone execution.

Two dependency conflicts were resolved during deployment. Installing
\texttt{ultralytics} brings in NumPy~2.x, against which \texttt{cv\_bridge} --- a
compiled extension built for the 1.x C ABI --- fails. Installing an older OpenCV
to compensate reintroduces NumPy~2 as a dependency, and OpenCV~5.0 separately
breaks \texttt{cv\_bridge} by renaming the type constants it uses to build its
conversion table. Versions were pinned to \texttt{numpy==1.26.4} and
\texttt{opencv-python==4.10.0.84}, the only combination satisfying both.

The model was additionally exported to NCNN for ARM-optimised inference. Because
NCNN fixes the input tensor shape at export time, the export geometry was
specified explicitly as \(544 \times 960\) to match the letterboxed frame; a
default export produces a square input in which approximately \(43\%\) of the
computation is spent on padding, at \(19.9\)\,GFLOPS against \(11.3\)\,GFLOPS for
the correct shape. Half-precision export was avoided, the Cortex-A72 lacking fp16
SIMD support.

% ---------------------------------------------------------------------
\subsection{Automation}
\label{sec:hitl-automation}

Bring-up was consolidated into a single script that starts the namespaces, ns-3,
Gazebo, SITL, the micro-ROS agent and the ground station receiver, creates the
VLANs and bridge, and launches the detector on the Pi over SSH. Pre-flight checks
were added for each failure encountered during development: presence of the host
address before Gazebo starts, since Fast DDS enumerates interfaces only at
participant creation and cannot use one assigned later; reachability of the Pi;
adequate socket buffer limits on both machines; and confirmation after startup
that Gazebo bound the intended interface.

% ---------------------------------------------------------------------
\subsection{Results}
\label{sec:hitl-results}

Detections were verified visually, by inspecting annotated frames, before any
quantitative result was accepted. An earlier configuration had reported up to
sixteen detections against five subjects; these were false positives on building
texture, eliminated by raising the confidence threshold and narrowing the camera
field of view.

\paragraph{Channel characterisation.}
Table~\ref{tab:hitl-links} compares the two paths, each measured over thirty ICMP
samples. Short samples are unreliable: the first packet includes address
resolution across the simulated channel, and a two-sample measurement of the same
configuration returned a mean of \(167\)\,ms against \(59\)\,ms for thirty.

\begin{table}[htbp]
\centering
\caption{Measured characteristics of the two links, thirty samples each.}
\label{tab:hitl-links}
\begin{tabular}{lrrrr}
\toprule
Path & Min (ms) & Mean (ms) & Max (ms) & Jitter (ms) \\
\midrule
Sensor, VLAN 10 (wired)   & 0.4  & 1.5  & 2.7   & 0.4  \\
Radio, VLAN 42 (via ns-3) & 11.7 & 77.4 & 242.6 & 54.8 \\
\bottomrule
\end{tabular}
\end{table}

Packet loss was zero on both paths, the aircraft remaining close to the ground
station throughout. The radio path is approximately \(53\times\) slower than the
wired path on average and varies approximately \(140\times\) more.

The distribution within the radio path is itself informative. Its minimum of
\(11.7\)\,ms against a mean of \(77.4\)\,ms is a spread of some \(6.6\times\) on
one link, where a channel applying a fixed delay would produce a narrow one.
The variation indicates that ns-3 is modelling contention, in which packets
queue and back off before transmission, rather than simply adding a constant
offset to each packet.

\paragraph{Sensor transport.}
With the camera at \(20\)\,Hz the wired path carried \(481\)\,Mbps at
\(42{,}544\) packets per second with zero dropped packets, occupying
approximately \(58\%\) of measured link capacity.

\paragraph{Inference latency.}
Table~\ref{tab:hitl-inference} reports latency for each combination of back-end
and input resolution. Changing back-end at constant resolution gives
approximately \(2.5\times\); halving the pixel count at constant back-end gives
approximately \(2.3\times\). Latency therefore scales close to linearly with
pixel count, a pixel ratio of \(0.47\) producing a time ratio of \(0.43\).

\begin{table}[htbp]
\centering
\caption{Inference latency on the Raspberry Pi 4B.}
\label{tab:hitl-inference}
\begin{tabular}{llrr}
\toprule
Back-end & Input & Latency (ms) & Rate (fps) \\
\midrule
PyTorch & $960 \times 544$ & 2540 & 0.39 \\
PyTorch & $640 \times 384$ & 1100 & 0.91 \\
NCNN    & $960 \times 544$ & 1000 & 1.00 \\
NCNN    & $640 \times 384$ & \multicolumn{2}{l}{not measured (est.\ 470)} \\
\bottomrule
\end{tabular}
\end{table}

\paragraph{Detection performance.}
Between two and four of the five subjects were detected per frame, a recall of
\(40\)--\(80\%\). The shortfall is attributable to scene geometry rather than
detector deficiency: the models are separated by approximately one metre, so from
the oblique camera view the nearer figures partially occlude those behind them,
and partially occluded subjects at approximately \(53\) pixels in height fall
below the confidence threshold. Because the number and position of subjects are
known exactly from the world definition, recall is a directly measurable quantity
rather than an estimate.

\paragraph{Transmitted volume.}
Each detection message is approximately \(118\)\,B, against \(2.76\)\,MB for the
corresponding frame, a reduction of approximately \(23{,}000\times\) in data
crossing the radio.

% ---------------------------------------------------------------------
\subsection{Contention between sensor transport and inference}
\label{sec:hitl-contention}

The two workloads on the edge node proved not to be independent. The same model,
on the same input, required \(623\)\,ms executed on a static image but
\(1{,}343\)\,ms within the live pipeline: a difference of \(720\)\,ms, or
approximately half the available inference throughput.

The edge node received frames at \(19.8\)\,Hz while completing inference at
approximately \(0.74\)\,Hz, discarding roughly \(96\%\) of what it received.
Discarded frames are not free: each requires reassembly of some \(2{,}150\) UDP
fragments, deserialisation of \(2.76\)\,MB and a further \(2.76\)\,MB
colour-space conversion, all competing with inference for the same four
processor cores.

Reducing the camera rate at source from \(20\)\,Hz to \(5\)\,Hz --- still
approximately seven times the rate the detector can consume --- reduced inference
latency to approximately \(1{,}000\)\,ms, recovering about a quarter of the
deficit, and reduced link occupancy from \(481\)\,Mbps to \(111\)\,Mbps. The
application-level frame-rate control provided by the relay node cannot achieve
this, as it discards frames after reception, by which point the transport and
deserialisation cost has already been incurred.

This interaction is not observable in pure simulation, where sensor data is
delivered through shared memory at negligible cost, nor in an isolated benchmark
of the detector, since the interference originates outside it. It appears only
when a physical processor performs both reception and inference concurrently.

% ---------------------------------------------------------------------
\subsection{Limitations}
\label{sec:hitl-limitations}

Clock synchronisation between the two machines has not yet been established.
End-to-end latency is computed from timestamps taken on different hosts, and
without synchronisation those values are not meaningful while remaining plausible
in magnitude; synchronisation is required before end-to-end latency is reported.

The edge configuration is characterised; the corresponding ground configuration,
in which compressed imagery crosses the simulated channel and inference is
performed at the ground station, remains to be measured.

Measurements were taken with the aircraft close to the ground station, giving
zero packet loss throughout. Behaviour as range increases and buildings obstruct
the path has not been characterised. The testbed supports a single physical edge
node; multi-node operation has not been evaluated.

%
%  HEADING LEVELS USED
%      \section     - the HITL section itself
%      \subsection  - its parts
%      \paragraph   - short findings within Results
%
%  REQUIRED PACKAGES (add to the main preamble if not already present)
%      \usepackage{booktabs}
%      \usepackage{graphicx}
%      \usepackage{tikz}
%      \usetikzlibrary{arrows.meta, positioning, fit, backgrounds}
%
%  All labels are prefixed hitl / tab:hitl- / fig:hitl- so they cannot
%  collide with labels elsewhere in the shared report.
% =====================================================================

\section{Hardware-in-the-Loop Edge Processing}
\label{sec:hitl}

The vision workload was moved off the simulation host and onto physical
hardware, so that the cost of onboard detection could be measured on a processor
representative of one that would actually fly. A Raspberry~Pi~4B performs
detection on imagery produced by the simulation, and returns its results across a
wireless channel simulated in ns-3.

% ---------------------------------------------------------------------
\subsection{Extending the testbed with physical hardware}
\label{sec:hitl-testbed}

The existing stack ran entirely on one machine: Gazebo for the environment and
sensors, ArduPilot SITL for the autopilot, ns-3 for the wireless channel, and
Linux network namespaces isolating the ground station from each aircraft. The
ns-3 scenario defined four nodes --- a ground control station and three aircraft
--- of which the second and third were unused placeholders.

A Raspberry~Pi~4B was introduced as the physical edge node and assigned to the
second aircraft slot. Its TAP device, previously created only to satisfy the
scenario's argument list, was brought into service as the attachment point. The
autopilot, physics and channel simulation remained on the host; only the vision
computation moved onto hardware.

The Pi was configured with Ubuntu~22.04 and ROS\,2 Humble to match the host
exactly, giving both machines the same DDS implementation. The host has no
Ethernet port, so a USB~3.0 gigabit adapter was fitted. An initial low-cost
adapter was replaced after \texttt{iperf3} measured \(4.5\)\,Mbps with \(24\%\)
loss, despite \texttt{ethtool} reporting \(100\)\,Mbps full duplex; the
replacement sustained \(834\)\,Mbps with no retransmissions.

% ---------------------------------------------------------------------
\subsection{Constructing the two-link network}
\label{sec:hitl-network}

Two logically separate paths were required between host and Pi. One carries
camera imagery and must remain unimpaired, since on a real aircraft it is a short
cable inside a single airframe rather than a radio. The other carries detection
results and must pass through the simulated channel, as that is the link the
experiment exists to characterise. The same reasoning applies to the
flight-dynamics traffic between SITL and Gazebo, which also bypasses ns-3.

The Pi provides a single Ethernet interface, so both paths were carried over one
cable using IEEE~802.1Q VLAN tagging. On the host, two virtual interfaces were
created on the physical adapter: \texttt{eth-cam}, tagged VLAN~10 and addressed
\(10.0.0.1\); and \texttt{eth-rf}, tagged VLAN~42, left without an address and
made a member of a new bridge, \texttt{br-uav2}. That bridge also received
\texttt{tap-uav2}, the ns-3 interface for the second aircraft node. On the Pi,
matching interfaces \texttt{eth0.10} (\(10.0.0.2\)) and \texttt{eth0.42}
(\(10.42.0.12\)) were created.

Attaching the physical interface to the bridge directly was rejected during
design. A Linux bridge floods frames to all member ports, so camera traffic would
have been delivered into the ns-3 TAP device as well, injecting approximately
\(111\)\,Mbps of imagery into the simulated channel.

The resulting arrangement enforces separation through routing rather than
filtering. Gazebo is addressed only on \(10.0.0.1\) and can reach the Pi only at
\(10.0.0.2\); the ground station namespace is addressed only on \(10.42.0.10\)
and can reach it only at \(10.42.0.12\). Neither can use the other's path.

\begin{figure}[htbp]
\centering
\begin{tikzpicture}[
  font=\small,
  box/.style={draw, rounded corners=2pt, align=center, inner sep=5pt},
  proc/.style={box, fill=black!4},
  net/.style={box, fill=black!10},
  arr/.style={-{Latex[length=2mm]}, thick},
]
\node[proc] (gz)   {Gazebo\\camera + physics};
\node[proc, below=14mm of gz] (gcs) {Ground station\\(namespace)\\\texttt{10.42.0.10}};
\node[net,  right=16mm of gcs] (ns3) {ns-3\\channel};
\node[net,  above=8mm of ns3]  (br)  {\texttt{br-uav2}\\+ TAP};
\node[proc, right=46mm of gz] (pi) {Raspberry Pi 4B\\YOLOv8n detector};

\draw[arr] (gz) -- node[above, align=center, font=\scriptsize]
  {VLAN 10 \quad camera, 2.76\,MB/frame\\unimpaired} (pi);
\draw[arr] (pi.south) |- node[pos=0.75, above, font=\scriptsize]
  {VLAN 42 \quad detections, 118\,B} (br.east);
\draw[arr] (br) -- (ns3);
\draw[arr] (ns3) -- (gcs);

\begin{scope}[on background layer]
  \node[draw, dashed, rounded corners, fit=(gz)(gcs)(ns3)(br),
        inner sep=7pt, label={[font=\scriptsize\itshape]above left:Host PC}] {};
\end{scope}
\end{tikzpicture}
\caption{The two-link arrangement. Camera imagery reaches the edge node over an
unimpaired wired path; detection results traverse the simulated channel. Both
share one physical cable, separated by VLAN tagging.}
\label{fig:hitl-architecture}
\end{figure}

Addressing was made persistent using NetworkManager on the host and netplan on
the Pi, the two distributions using different network managers. Addresses
assigned imperatively were found to be removed within minutes, as each manager
reconciles kernel state against its own configuration and deletes entries it did
not create.

% ---------------------------------------------------------------------
\subsection{Configuring the middleware}
\label{sec:hitl-dds}

A Fast DDS profile was written and applied on both machines. It pins participants
to the intended interfaces through an interface whitelist, since both machines
are attached to WiFi as well as the wired link and DDS otherwise advertises on
every available interface. The profile also disables the shared-memory transport,
which is not isolated by network namespaces and would allow participants to
exchange data without traversing the simulated channel.

Kernel socket buffers were raised on both machines from the \(208\)\,kB default
to \(512\)\,MB. A single \(1280 \times 720\) frame is \(2.76\)\,MB and fragments
into approximately \(1{,}900\) UDP datagrams, of which the default buffer holds
about \(7\%\); reassembly never completed and subscribers received nothing, with
no error reported.

The camera subscription on the Pi was changed from the ROS\,2 default of
\texttt{RELIABLE} with queue depth ten to \texttt{BEST\_EFFORT} with depth one.
Under the default policy the publisher retains each sample until the subscriber
acknowledges it, and because the Pi acknowledges at approximately \(1\)\,Hz, the
camera output of the entire simulation was throttled to \(0.27\)\,Hz.
\texttt{RELIABLE} was retained for detection messages, which must not be lost.

% ---------------------------------------------------------------------
\subsection{Deploying the detector}
\label{sec:hitl-detector}

The detection node runs YOLOv8n restricted to the COCO \emph{person} class,
subscribing to the camera topic and publishing results as JSON. It was deployed
to the Pi and instrumented to record inference duration, measured around the
model invocation alone so that figures remain comparable between back-ends and
against standalone execution.

Two dependency conflicts were resolved during deployment. Installing
\texttt{ultralytics} brings in NumPy~2.x, against which \texttt{cv\_bridge} --- a
compiled extension built for the 1.x C ABI --- fails. Installing an older OpenCV
to compensate reintroduces NumPy~2 as a dependency, and OpenCV~5.0 separately
breaks \texttt{cv\_bridge} by renaming the type constants it uses to build its
conversion table. Versions were pinned to \texttt{numpy==1.26.4} and
\texttt{opencv-python==4.10.0.84}, the only combination satisfying both.

The model was additionally exported to NCNN for ARM-optimised inference. Because
NCNN fixes the input tensor shape at export time, the export geometry was
specified explicitly as \(544 \times 960\) to match the letterboxed frame; a
default export produces a square input in which approximately \(43\%\) of the
computation is spent on padding, at \(19.9\)\,GFLOPS against \(11.3\)\,GFLOPS for
the correct shape. Half-precision export was avoided, the Cortex-A72 lacking fp16
SIMD support.

% ---------------------------------------------------------------------
\subsection{Automation}
\label{sec:hitl-automation}

Bring-up was consolidated into a single script that starts the namespaces, ns-3,
Gazebo, SITL, the micro-ROS agent and the ground station receiver, creates the
VLANs and bridge, and launches the detector on the Pi over SSH. Pre-flight checks
were added for each failure encountered during development: presence of the host
address before Gazebo starts, since Fast DDS enumerates interfaces only at
participant creation and cannot use one assigned later; reachability of the Pi;
adequate socket buffer limits on both machines; and confirmation after startup
that Gazebo bound the intended interface.

% ---------------------------------------------------------------------
\subsection{Results}
\label{sec:hitl-results}

Detections were verified visually, by inspecting annotated frames, before any
quantitative result was accepted. An earlier configuration had reported up to
sixteen detections against five subjects; these were false positives on building
texture, eliminated by raising the confidence threshold and narrowing the camera
field of view.

\paragraph{Channel characterisation.}
Table~\ref{tab:hitl-links} compares the two paths, each measured over thirty ICMP
samples. Short samples are unreliable: the first packet includes address
resolution across the simulated channel, and a two-sample measurement of the same
configuration returned a mean of \(167\)\,ms against \(59\)\,ms for thirty.

\begin{table}[htbp]
\centering
\caption{Measured characteristics of the two links, thirty samples each.}
\label{tab:hitl-links}
\begin{tabular}{lrrrr}
\toprule
Path & Min (ms) & Mean (ms) & Max (ms) & Jitter (ms) \\
\midrule
Sensor, VLAN 10 (wired)   & 0.4  & 1.5  & 2.7   & 0.4  \\
Radio, VLAN 42 (via ns-3) & 11.7 & 77.4 & 242.6 & 54.8 \\
\bottomrule
\end{tabular}
\end{table}

Packet loss was zero on both paths, the aircraft remaining close to the ground
station throughout. The radio path is approximately \(53\times\) slower than the
wired path on average and varies approximately \(140\times\) more.

The distribution within the radio path is itself informative. Its minimum of
\(11.7\)\,ms against a mean of \(77.4\)\,ms is a spread of some \(6.6\times\) on
one link, where a channel applying a fixed delay would produce a narrow one.
The variation indicates that ns-3 is modelling contention, in which packets
queue and back off before transmission, rather than simply adding a constant
offset to each packet.

\paragraph{Sensor transport.}
With the camera at \(20\)\,Hz the wired path carried \(481\)\,Mbps at
\(42{,}544\) packets per second with zero dropped packets, occupying
approximately \(58\%\) of measured link capacity.

\paragraph{Inference latency.}
Table~\ref{tab:hitl-inference} reports latency for each combination of back-end
and input resolution. Changing back-end at constant resolution gives
approximately \(2.5\times\); halving the pixel count at constant back-end gives
approximately \(2.3\times\). Latency therefore scales close to linearly with
pixel count, a pixel ratio of \(0.47\) producing a time ratio of \(0.43\).

\begin{table}[htbp]
\centering
\caption{Inference latency on the Raspberry Pi 4B.}
\label{tab:hitl-inference}
\begin{tabular}{llrr}
\toprule
Back-end & Input & Latency (ms) & Rate (fps) \\
\midrule
PyTorch & $960 \times 544$ & 2540 & 0.39 \\
PyTorch & $640 \times 384$ & 1100 & 0.91 \\
NCNN    & $960 \times 544$ & 1000 & 1.00 \\
NCNN    & $640 \times 384$ & \multicolumn{2}{l}{not measured (est.\ 470)} \\
\bottomrule
\end{tabular}
\end{table}

\paragraph{Detection performance.}
Between two and four of the five subjects were detected per frame, a recall of
\(40\)--\(80\%\). The shortfall is attributable to scene geometry rather than
detector deficiency: the models are separated by approximately one metre, so from
the oblique camera view the nearer figures partially occlude those behind them,
and partially occluded subjects at approximately \(53\) pixels in height fall
below the confidence threshold. Because the number and position of subjects are
known exactly from the world definition, recall is a directly measurable quantity
rather than an estimate.

\paragraph{Transmitted volume.}
Each detection message is approximately \(118\)\,B, against \(2.76\)\,MB for the
corresponding frame, a reduction of approximately \(23{,}000\times\) in data
crossing the radio.

% ---------------------------------------------------------------------
\subsection{Contention between sensor transport and inference}
\label{sec:hitl-contention}

The two workloads on the edge node proved not to be independent. The same model,
on the same input, required \(623\)\,ms executed on a static image but
\(1{,}343\)\,ms within the live pipeline: a difference of \(720\)\,ms, or
approximately half the available inference throughput.

The edge node received frames at \(19.8\)\,Hz while completing inference at
approximately \(0.74\)\,Hz, discarding roughly \(96\%\) of what it received.
Discarded frames are not free: each requires reassembly of some \(2{,}150\) UDP
fragments, deserialisation of \(2.76\)\,MB and a further \(2.76\)\,MB
colour-space conversion, all competing with inference for the same four
processor cores.

Reducing the camera rate at source from \(20\)\,Hz to \(5\)\,Hz --- still
approximately seven times the rate the detector can consume --- reduced inference
latency to approximately \(1{,}000\)\,ms, recovering about a quarter of the
deficit, and reduced link occupancy from \(481\)\,Mbps to \(111\)\,Mbps. The
application-level frame-rate control provided by the relay node cannot achieve
this, as it discards frames after reception, by which point the transport and
deserialisation cost has already been incurred.

This interaction is not observable in pure simulation, where sensor data is
delivered through shared memory at negligible cost, nor in an isolated benchmark
of the detector, since the interference originates outside it. It appears only
when a physical processor performs both reception and inference concurrently.

% ---------------------------------------------------------------------
\subsection{Limitations}
\label{sec:hitl-limitations}

Clock synchronisation between the two machines has not yet been established.
End-to-end latency is computed from timestamps taken on different hosts, and
without synchronisation those values are not meaningful while remaining plausible
in magnitude; synchronisation is required before end-to-end latency is reported.

The edge configuration is characterised; the corresponding ground configuration,
in which compressed imagery crosses the simulated channel and inference is
performed at the ground station, remains to be measured.

Measurements were taken with the aircraft close to the ground station, giving
zero packet loss throughout. Behaviour as range increases and buildings obstruct
the path has not been characterised. The testbed supports a single physical edge
node; multi-node operation has not been evaluated.

%
%  HEADING LEVELS USED
%      \section     - the HITL section itself
%      \subsection  - its parts
%      \paragraph   - short findings within Results
%
%  REQUIRED PACKAGES (add to the main preamble if not already present)
%      \usepackage{booktabs}
%      \usepackage{graphicx}
%      \usepackage{tikz}
%      \usetikzlibrary{arrows.meta, positioning, fit, backgrounds}
%
%  All labels are prefixed hitl / tab:hitl- / fig:hitl- so they cannot
%  collide with labels elsewhere in the shared report.
% =====================================================================

\section{Hardware-in-the-Loop Edge Processing}
\label{sec:hitl}

The vision workload was moved off the simulation host and onto physical
hardware, so that the cost of onboard detection could be measured on a processor
representative of one that would actually fly. A Raspberry~Pi~4B performs
detection on imagery produced by the simulation, and returns its results across a
wireless channel simulated in ns-3.

% ---------------------------------------------------------------------
\subsection{Extending the testbed with physical hardware}
\label{sec:hitl-testbed}

The existing stack ran entirely on one machine: Gazebo for the environment and
sensors, ArduPilot SITL for the autopilot, ns-3 for the wireless channel, and
Linux network namespaces isolating the ground station from each aircraft. The
ns-3 scenario defined four nodes --- a ground control station and three aircraft
--- of which the second and third were unused placeholders.

A Raspberry~Pi~4B was introduced as the physical edge node and assigned to the
second aircraft slot. Its TAP device, previously created only to satisfy the
scenario's argument list, was brought into service as the attachment point. The
autopilot, physics and channel simulation remained on the host; only the vision
computation moved onto hardware.

The Pi was configured with Ubuntu~22.04 and ROS\,2 Humble to match the host
exactly, giving both machines the same DDS implementation. The host has no
Ethernet port, so a USB~3.0 gigabit adapter was fitted. An initial low-cost
adapter was replaced after \texttt{iperf3} measured \(4.5\)\,Mbps with \(24\%\)
loss, despite \texttt{ethtool} reporting \(100\)\,Mbps full duplex; the
replacement sustained \(834\)\,Mbps with no retransmissions.

% ---------------------------------------------------------------------
\subsection{Constructing the two-link network}
\label{sec:hitl-network}

Two logically separate paths were required between host and Pi. One carries
camera imagery and must remain unimpaired, since on a real aircraft it is a short
cable inside a single airframe rather than a radio. The other carries detection
results and must pass through the simulated channel, as that is the link the
experiment exists to characterise. The same reasoning applies to the
flight-dynamics traffic between SITL and Gazebo, which also bypasses ns-3.

The Pi provides a single Ethernet interface, so both paths were carried over one
cable using IEEE~802.1Q VLAN tagging. On the host, two virtual interfaces were
created on the physical adapter: \texttt{eth-cam}, tagged VLAN~10 and addressed
\(10.0.0.1\); and \texttt{eth-rf}, tagged VLAN~42, left without an address and
made a member of a new bridge, \texttt{br-uav2}. That bridge also received
\texttt{tap-uav2}, the ns-3 interface for the second aircraft node. On the Pi,
matching interfaces \texttt{eth0.10} (\(10.0.0.2\)) and \texttt{eth0.42}
(\(10.42.0.12\)) were created.

Attaching the physical interface to the bridge directly was rejected during
design. A Linux bridge floods frames to all member ports, so camera traffic would
have been delivered into the ns-3 TAP device as well, injecting approximately
\(111\)\,Mbps of imagery into the simulated channel.

The resulting arrangement enforces separation through routing rather than
filtering. Gazebo is addressed only on \(10.0.0.1\) and can reach the Pi only at
\(10.0.0.2\); the ground station namespace is addressed only on \(10.42.0.10\)
and can reach it only at \(10.42.0.12\). Neither can use the other's path.

\begin{figure}[htbp]
\centering
\begin{tikzpicture}[
  font=\small,
  box/.style={draw, rounded corners=2pt, align=center, inner sep=5pt},
  proc/.style={box, fill=black!4},
  net/.style={box, fill=black!10},
  arr/.style={-{Latex[length=2mm]}, thick},
]
\node[proc] (gz)   {Gazebo\\camera + physics};
\node[proc, below=14mm of gz] (gcs) {Ground station\\(namespace)\\\texttt{10.42.0.10}};
\node[net,  right=16mm of gcs] (ns3) {ns-3\\channel};
\node[net,  above=8mm of ns3]  (br)  {\texttt{br-uav2}\\+ TAP};
\node[proc, right=46mm of gz] (pi) {Raspberry Pi 4B\\YOLOv8n detector};

\draw[arr] (gz) -- node[above, align=center, font=\scriptsize]
  {VLAN 10 \quad camera, 2.76\,MB/frame\\unimpaired} (pi);
\draw[arr] (pi.south) |- node[pos=0.75, above, font=\scriptsize]
  {VLAN 42 \quad detections, 118\,B} (br.east);
\draw[arr] (br) -- (ns3);
\draw[arr] (ns3) -- (gcs);

\begin{scope}[on background layer]
  \node[draw, dashed, rounded corners, fit=(gz)(gcs)(ns3)(br),
        inner sep=7pt, label={[font=\scriptsize\itshape]above left:Host PC}] {};
\end{scope}
\end{tikzpicture}
\caption{The two-link arrangement. Camera imagery reaches the edge node over an
unimpaired wired path; detection results traverse the simulated channel. Both
share one physical cable, separated by VLAN tagging.}
\label{fig:hitl-architecture}
\end{figure}

Addressing was made persistent using NetworkManager on the host and netplan on
the Pi, the two distributions using different network managers. Addresses
assigned imperatively were found to be removed within minutes, as each manager
reconciles kernel state against its own configuration and deletes entries it did
not create.

% ---------------------------------------------------------------------
\subsection{Configuring the middleware}
\label{sec:hitl-dds}

A Fast DDS profile was written and applied on both machines. It pins participants
to the intended interfaces through an interface whitelist, since both machines
are attached to WiFi as well as the wired link and DDS otherwise advertises on
every available interface. The profile also disables the shared-memory transport,
which is not isolated by network namespaces and would allow participants to
exchange data without traversing the simulated channel.

Kernel socket buffers were raised on both machines from the \(208\)\,kB default
to \(512\)\,MB. A single \(1280 \times 720\) frame is \(2.76\)\,MB and fragments
into approximately \(1{,}900\) UDP datagrams, of which the default buffer holds
about \(7\%\); reassembly never completed and subscribers received nothing, with
no error reported.

The camera subscription on the Pi was changed from the ROS\,2 default of
\texttt{RELIABLE} with queue depth ten to \texttt{BEST\_EFFORT} with depth one.
Under the default policy the publisher retains each sample until the subscriber
acknowledges it, and because the Pi acknowledges at approximately \(1\)\,Hz, the
camera output of the entire simulation was throttled to \(0.27\)\,Hz.
\texttt{RELIABLE} was retained for detection messages, which must not be lost.

% ---------------------------------------------------------------------
\subsection{Deploying the detector}
\label{sec:hitl-detector}

The detection node runs YOLOv8n restricted to the COCO \emph{person} class,
subscribing to the camera topic and publishing results as JSON. It was deployed
to the Pi and instrumented to record inference duration, measured around the
model invocation alone so that figures remain comparable between back-ends and
against standalone execution.

Two dependency conflicts were resolved during deployment. Installing
\texttt{ultralytics} brings in NumPy~2.x, against which \texttt{cv\_bridge} --- a
compiled extension built for the 1.x C ABI --- fails. Installing an older OpenCV
to compensate reintroduces NumPy~2 as a dependency, and OpenCV~5.0 separately
breaks \texttt{cv\_bridge} by renaming the type constants it uses to build its
conversion table. Versions were pinned to \texttt{numpy==1.26.4} and
\texttt{opencv-python==4.10.0.84}, the only combination satisfying both.

The model was additionally exported to NCNN for ARM-optimised inference. Because
NCNN fixes the input tensor shape at export time, the export geometry was
specified explicitly as \(544 \times 960\) to match the letterboxed frame; a
default export produces a square input in which approximately \(43\%\) of the
computation is spent on padding, at \(19.9\)\,GFLOPS against \(11.3\)\,GFLOPS for
the correct shape. Half-precision export was avoided, the Cortex-A72 lacking fp16
SIMD support.

% ---------------------------------------------------------------------
\subsection{Automation}
\label{sec:hitl-automation}

Bring-up was consolidated into a single script that starts the namespaces, ns-3,
Gazebo, SITL, the micro-ROS agent and the ground station receiver, creates the
VLANs and bridge, and launches the detector on the Pi over SSH. Pre-flight checks
were added for each failure encountered during development: presence of the host
address before Gazebo starts, since Fast DDS enumerates interfaces only at
participant creation and cannot use one assigned later; reachability of the Pi;
adequate socket buffer limits on both machines; and confirmation after startup
that Gazebo bound the intended interface.

% ---------------------------------------------------------------------
\subsection{Results}
\label{sec:hitl-results}

Detections were verified visually, by inspecting annotated frames, before any
quantitative result was accepted. An earlier configuration had reported up to
sixteen detections against five subjects; these were false positives on building
texture, eliminated by raising the confidence threshold and narrowing the camera
field of view.

\paragraph{Channel characterisation.}
Table~\ref{tab:hitl-links} compares the two paths, each measured over thirty ICMP
samples. Short samples are unreliable: the first packet includes address
resolution across the simulated channel, and a two-sample measurement of the same
configuration returned a mean of \(167\)\,ms against \(59\)\,ms for thirty.

\begin{table}[htbp]
\centering
\caption{Measured characteristics of the two links, thirty samples each.}
\label{tab:hitl-links}
\begin{tabular}{lrrrr}
\toprule
Path & Min (ms) & Mean (ms) & Max (ms) & Jitter (ms) \\
\midrule
Sensor, VLAN 10 (wired)   & 0.4  & 1.5  & 2.7   & 0.4  \\
Radio, VLAN 42 (via ns-3) & 11.7 & 77.4 & 242.6 & 54.8 \\
\bottomrule
\end{tabular}
\end{table}

Packet loss was zero on both paths, the aircraft remaining close to the ground
station throughout. The radio path is approximately \(53\times\) slower than the
wired path on average and varies approximately \(140\times\) more.

The distribution within the radio path is itself informative. Its minimum of
\(11.7\)\,ms against a mean of \(77.4\)\,ms is a spread of some \(6.6\times\) on
one link, where a channel applying a fixed delay would produce a narrow one.
The variation indicates that ns-3 is modelling contention, in which packets
queue and back off before transmission, rather than simply adding a constant
offset to each packet.

\paragraph{Sensor transport.}
With the camera at \(20\)\,Hz the wired path carried \(481\)\,Mbps at
\(42{,}544\) packets per second with zero dropped packets, occupying
approximately \(58\%\) of measured link capacity.

\paragraph{Inference latency.}
Table~\ref{tab:hitl-inference} reports latency for each combination of back-end
and input resolution. Changing back-end at constant resolution gives
approximately \(2.5\times\); halving the pixel count at constant back-end gives
approximately \(2.3\times\). Latency therefore scales close to linearly with
pixel count, a pixel ratio of \(0.47\) producing a time ratio of \(0.43\).

\begin{table}[htbp]
\centering
\caption{Inference latency on the Raspberry Pi 4B.}
\label{tab:hitl-inference}
\begin{tabular}{llrr}
\toprule
Back-end & Input & Latency (ms) & Rate (fps) \\
\midrule
PyTorch & $960 \times 544$ & 2540 & 0.39 \\
PyTorch & $640 \times 384$ & 1100 & 0.91 \\
NCNN    & $960 \times 544$ & 1000 & 1.00 \\
NCNN    & $640 \times 384$ & \multicolumn{2}{l}{not measured (est.\ 470)} \\
\bottomrule
\end{tabular}
\end{table}

\paragraph{Detection performance.}
Between two and four of the five subjects were detected per frame, a recall of
\(40\)--\(80\%\). The shortfall is attributable to scene geometry rather than
detector deficiency: the models are separated by approximately one metre, so from
the oblique camera view the nearer figures partially occlude those behind them,
and partially occluded subjects at approximately \(53\) pixels in height fall
below the confidence threshold. Because the number and position of subjects are
known exactly from the world definition, recall is a directly measurable quantity
rather than an estimate.

\paragraph{Transmitted volume.}
Each detection message is approximately \(118\)\,B, against \(2.76\)\,MB for the
corresponding frame, a reduction of approximately \(23{,}000\times\) in data
crossing the radio.

% ---------------------------------------------------------------------
\subsection{Contention between sensor transport and inference}
\label{sec:hitl-contention}

The two workloads on the edge node proved not to be independent. The same model,
on the same input, required \(623\)\,ms executed on a static image but
\(1{,}343\)\,ms within the live pipeline: a difference of \(720\)\,ms, or
approximately half the available inference throughput.

The edge node received frames at \(19.8\)\,Hz while completing inference at
approximately \(0.74\)\,Hz, discarding roughly \(96\%\) of what it received.
Discarded frames are not free: each requires reassembly of some \(2{,}150\) UDP
fragments, deserialisation of \(2.76\)\,MB and a further \(2.76\)\,MB
colour-space conversion, all competing with inference for the same four
processor cores.

Reducing the camera rate at source from \(20\)\,Hz to \(5\)\,Hz --- still
approximately seven times the rate the detector can consume --- reduced inference
latency to approximately \(1{,}000\)\,ms, recovering about a quarter of the
deficit, and reduced link occupancy from \(481\)\,Mbps to \(111\)\,Mbps. The
application-level frame-rate control provided by the relay node cannot achieve
this, as it discards frames after reception, by which point the transport and
deserialisation cost has already been incurred.

This interaction is not observable in pure simulation, where sensor data is
delivered through shared memory at negligible cost, nor in an isolated benchmark
of the detector, since the interference originates outside it. It appears only
when a physical processor performs both reception and inference concurrently.

% ---------------------------------------------------------------------
\subsection{Limitations}
\label{sec:hitl-limitations}

Clock synchronisation between the two machines has not yet been established.
End-to-end latency is computed from timestamps taken on different hosts, and
without synchronisation those values are not meaningful while remaining plausible
in magnitude; synchronisation is required before end-to-end latency is reported.

The edge configuration is characterised; the corresponding ground configuration,
in which compressed imagery crosses the simulated channel and inference is
performed at the ground station, remains to be measured.

Measurements were taken with the aircraft close to the ground station, giving
zero packet loss throughout. Behaviour as range increases and buildings obstruct
the path has not been characterised. The testbed supports a single physical edge
node; multi-node operation has not been evaluated.
; it cross-references
%  \ref{sec:hitl} and its tables. If the two are placed in different
%  chapters the \ref commands still resolve, LaTeX just needs two passes.
%
%  REQUIRED PACKAGES (add to the main preamble if not already present)
%      \usepackage{booktabs}
%      \usepackage{array}
%      \usepackage{graphicx}
%      \usepackage{subcaption}          % only for the side-by-side figure
%
%  FIGURES
%  Place image files in  docs/report/figures/  and add to the preamble:
%      \graphicspath{{figures/}}
%  Each \includegraphics below names the file it expects. See the checklist
%  at the end of this file for what each image should show.
%
%  All labels are prefixed test / tab:test- so they cannot collide with
%  labels elsewhere in the shared report.
% =====================================================================

\newcommand{\pass}{\textsc{pass}}
\newcommand{\fail}{\textsc{fail}}

\section{Testing, Experimentation and Validation}
\label{sec:test}

This section sets out how the hardware-in-the-loop testbed described in
Section~\ref{sec:hitl} was verified. It distinguishes three activities.
\emph{Testing} establishes that each component behaves as specified.
\emph{Experimentation} varies a controlled parameter to obtain the
measurements the project set out to produce. \emph{Validation} establishes
that those measurements mean what they are claimed to mean, which in a
partly simulated system cannot be taken for granted.

% ---------------------------------------------------------------------
\subsection{Testing strategy}
\label{sec:test-strategy}

Testing was performed bottom-up, in five layers, each of which had to pass
before the next was attempted. The order is dictated by the failure modes
observed during development rather than chosen for tidiness: in this system
a fault at a lower layer does not produce an error at the layer above, it
produces plausible but wrong behaviour there.

Three examples motivated the approach. When the kernel socket buffer was too
small, frame reassembly never completed and the subscriber received nothing
at all --- with no exception raised and no message logged. When the
subscriber's quality-of-service policy was left at the ROS\,2 default, the
publisher throttled itself and the whole simulation slowed, which presents
as a rendering problem rather than a network one. When the detector's
publisher and the ground station's subscriber were given mismatched topic
names, both processes ran without error indefinitely and simply never
exchanged a message. In each case the visible symptom appeared several
layers above the cause.

The layers are: the physical link (L1); logical separation of the two paths
(L2); middleware transport (L3); detector function (L4); and the end-to-end
path from camera to ground station (L5). Every test is recorded below with
the command used, so each is repeatable rather than merely asserted.

% ---------------------------------------------------------------------
\subsection{Test environment}
\label{sec:test-env}

Table~\ref{tab:test-env} records the configuration under which all results
in this report were obtained. Both machines run the same operating system
release and the same ROS\,2 distribution, so that the middleware
implementation and version are identical at both ends of the link and cannot
themselves be a source of difference.

\begin{table}[htbp]
\centering
\caption{Test environment.}
\label{tab:test-env}
\begin{tabular}{ll}
\toprule
Item & Configuration \\
\midrule
Simulation host    & x86-64 laptop, no onboard Ethernet \\
Edge node          & Raspberry Pi 4B, Cortex-A72 quad core \\
Link adapter       & USB 3.0 gigabit Ethernet, host side \\
Operating system   & Ubuntu 22.04 (both machines) \\
Middleware         & ROS\,2 Humble, Fast DDS 2.6 (both machines) \\
Simulator          & Gazebo Classic, lockstep with ArduPilot SITL \\
Channel simulator  & ns-3, TapBridge, 4-node scenario \\
Detector           & YOLOv8n, COCO class \emph{person} only \\
Camera             & $1280 \times 720$, 2.76\,MB per frame \\
Ground truth       & 5 human models, known world coordinates \\
\bottomrule
\end{tabular}
\end{table}

The Pi occupies the second aircraft node of the ns-3 scenario. Its single
Ethernet port is divided by 802.1Q tagging into \texttt{eth0.10}
($10.0.0.2$, camera, unimpaired) and \texttt{eth0.42} ($10.42.0.12$,
detections, through ns-3). The corresponding host interfaces are
\texttt{eth-cam} ($10.0.0.1$) and \texttt{eth-rf}, the latter bridged to
\texttt{tap-uav2}. Interfaces are created by
\texttt{scripts/netns/pi\_hitl\_link.sh} on the Pi and by step~1c of
\texttt{scripts/netns/launch\_single\_uav\_netns.sh} on the host.

\begin{figure}[htbp]
\centering
\includegraphics[width=0.7\textwidth]{testbed_hardware.jpg}
\caption{The physical testbed. The simulation host is connected to the
Raspberry~Pi~4B edge node by a single Ethernet cable through a USB~3.0
gigabit adapter. That one cable carries both the unimpaired sensor link
(VLAN~10) and the simulated radio link (VLAN~42).}
\label{fig:test-hardware}
\end{figure}

% ---------------------------------------------------------------------
\subsection{Layer 1 --- Physical link}
\label{sec:test-l1}

The first layer establishes that the cable between the two machines carries
data at the rate the experiment requires. This matters because the camera
stream alone occupies roughly half of a gigabit link; an adapter that
negotiates successfully but performs poorly would corrupt every result
downstream.

\begin{table}[htbp]
\centering
\caption{Layer 1 test cases --- physical link.}
\label{tab:test-l1}
\begin{tabular}{@{}lp{40mm}p{38mm}l@{}}
\toprule
ID & Test and command & Result & Verdict \\
\midrule
T1.1 & Link negotiation. \newline \texttt{ethtool <if>}
     & $1000$\,Mbps full duplex, link detected
     & \pass \\[2pt]
T1.2 & Throughput, first adapter. \newline \texttt{iperf3 -c 10.0.0.2}
     & $4.5$\,Mbps, $24\%$ loss, despite \texttt{ethtool} reporting
       $100$\,Mbps
     & \fail \\[2pt]
T1.3 & Throughput, replacement adapter. \newline \texttt{iperf3 -c 10.0.0.2}
     & $834$\,Mbps, no retransmissions
     & \pass \\[2pt]
T1.4 & Interface state after reboot. \newline \texttt{ip -br addr show}
     & Addresses present on both machines after persistence was configured
     & \pass \\
\bottomrule
\end{tabular}
\end{table}

Test~T1.2 is retained in this report because it is the clearest instance of
the reporting problem described above: \texttt{ethtool} reported a healthy
negotiated link while the adapter delivered $4.5$\,Mbps with $24\%$ loss.
Link negotiation and link performance are separate properties, and only the
second was fit for purpose. Had this been discovered later, the resulting
frame loss would most naturally have been attributed to the middleware or to
the detector.

Test~T1.4 was added after addresses assigned with \texttt{ip addr add} were
observed to disappear within minutes. Each machine runs a network manager
that reconciles kernel state against its own stored configuration ---
NetworkManager on the host, netplan on the Pi --- and removes addresses it
did not create. The test therefore checks persistence, not merely presence.

% ---------------------------------------------------------------------
\subsection{Layer 2 --- Logical separation of the two paths}
\label{sec:test-l2}

The central requirement of the experiment is that detection results cross
the simulated wireless channel while camera imagery does not. If camera
traffic were to leak into ns-3 the channel measurements would describe a
link carrying $481$\,Mbps of video, which is not the link the aircraft
would have. This layer therefore receives the most attention, and is tested
positively (the intended traffic is present) and negatively (the unintended
traffic is absent).

\begin{table}[htbp]
\centering
\caption{Layer 2 test cases --- path separation.}
\label{tab:test-l2}
\begin{tabular}{@{}lp{40mm}p{38mm}l@{}}
\toprule
ID & Test and command & Result & Verdict \\
\midrule
T2.1 & VLAN interfaces exist and are addressed. \newline
       \texttt{ip -br addr show}
     & \texttt{eth0.10} $10.0.0.2$, \texttt{eth0.42} $10.42.0.12$ on the Pi;
       \texttt{eth-cam} $10.0.0.1$ on the host
     & \pass \\[2pt]
T2.2 & Camera path reachable. \newline \texttt{ping -c 30 10.0.0.1}
     & Mean $1.5$\,ms, $0\%$ loss
     & \pass \\[2pt]
T2.3 & Radio path reachable through ns-3. \newline
       \texttt{ping -c 30 10.42.0.10}
     & Mean $77.4$\,ms, $0\%$ loss
     & \pass \\[2pt]
T2.4 & The two paths are distinguishable. \newline
       Compare T2.2 and T2.3
     & $53\times$ latency, $140\times$ jitter
     & \pass \\[2pt]
T2.5 & Camera traffic absent from the simulated channel. \newline
       \texttt{tcpdump -i tap-uav2 -nn}
     & No $10.0.0.0/24$ traffic observed on the ns-3 interface
     & \pass \\[2pt]
T2.6 & Volume crossing each path. \newline
       \texttt{ip -s link show eth-cam} \newline
       \texttt{ip -s link show tap-uav2}
     & Byte counters differ by approximately four orders of magnitude
     & \pass \\
\bottomrule
\end{tabular}
\end{table}

Test~T2.4 deserves comment, because two paths that both work are not by
themselves evidence of separation --- they must also be shown to differ in
the intended way. The measured contrast is large and in the expected
direction: the wired path is fast and stable, the simulated radio slow and
variable. Table~\ref{tab:hitl-links} gives the full figures.

Tests~T2.5 and T2.6 are the negative tests. Separation here is enforced by
addressing and bridge membership rather than by firewall rules, so it is
established by observing what crosses the interfaces rather than by
inspecting a rule set. \texttt{tap-uav2} is the point at which traffic
enters ns-3; if imagery is absent there, it is absent from the simulated
channel. T2.6 makes the same point quantitatively, since a leak of even one
frame per second would be immediately visible in the byte counters given
that a frame is $2.76$\,MB against $118$\,B for a detection message.

The alternative design --- attaching the physical interface to the bridge
directly --- was rejected on exactly this ground during design rather than
during testing. A Linux bridge floods frames to all member ports, so the
arrangement would have failed T2.5 by construction.

% ---------------------------------------------------------------------
\subsection{Layer 3 --- Middleware transport}
\label{sec:test-l3}

With the links established, this layer tests that DDS uses them as intended
and that large messages survive the crossing. All three tests here
correspond to faults that were encountered and produced no error message.

\begin{table}[htbp]
\centering
\caption{Layer 3 test cases --- middleware transport.}
\label{tab:test-l3}
\begin{tabular}{@{}lp{40mm}p{38mm}l@{}}
\toprule
ID & Test and command & Result & Verdict \\
\midrule
T3.1 & Receive buffers adequate for a $2.76$\,MB frame. \newline
       \texttt{sysctl net.core.rmem\_max}
     & Raised from $208$\,kB to $512$\,MB on both machines
     & \pass \\[2pt]
T3.2 & \emph{Send} buffers adequate. \newline
       \texttt{sysctl net.core.wmem\_max}
     & Initially $208$\,kB; raised to $512$\,MB on both machines
     & \fail\,$\rightarrow$\,\pass \\[2pt]
T3.3 & Publisher bound to the intended interface. \newline
       \texttt{ss -ulnp | grep gzserver}
     & \texttt{gzserver} bound to $10.0.0.1$
     & \pass \\[2pt]
T3.4 & Frames arrive complete and are not dropped. \newline
       \texttt{ip -s link show eth0.10}
     & Zero dropped, zero errors; $42{,}544$ packets\,s$^{-1}$ at
       $481$\,Mbps when the source ran at $20$\,Hz
     & \pass \\[2pt]
T3.5 & Delivered frame rate matches the source. \newline
       \texttt{ros2 topic hz /uav1/camera/image\_raw}
     & $5.000$\,Hz against a $5$\,Hz source, standard deviation
       $4.5$\,ms
     & \pass \\[2pt]
T3.6 & Subscriber QoS as specified. \newline
       \texttt{ros2 topic info -v <topic>}
     & \texttt{BEST\_EFFORT}, depth $1$ on the camera subscription
     & \pass \\[2pt]
T3.7 & Publisher and subscriber matched. \newline
       \texttt{ros2 topic info -v /detections/uav1}
     & One publisher, one subscriber
     & \pass \\
\bottomrule
\end{tabular}
\end{table}

Test~T3.3 exists because Fast DDS~2.6 enumerates network interfaces once,
at participant creation, and has no dynamic detection. If Gazebo starts
before $10.0.0.1$ has been assigned it can never bind that address, and the
Pi silently receives no imagery for the entire run. The test is therefore
also implemented as an automated pre-flight check in
\texttt{scripts/netns/run\_hitl.sh}, both before start-up, on the address,
and after start-up, on the actual socket binding.

Test~T3.5 is a rate check rather than a delivery check, and it was
introduced after the ROS\,2 default QoS was found to throttle the camera
output of the entire simulation to $0.27$\,Hz. That failure is invisible to
T3.4, since the packets that are sent do all arrive; only the rate reveals
it. A related observation gave the diagnosis: with the Pi disconnected, the
publisher continued to emit at intervals of $3.144$\,s with a standard
deviation of $0.4$\,ms. A spacing that regular is a retransmission timer
rather than a rendering rate, which located the fault in the middleware
rather than in Gazebo.

Test~T3.7 was added after the detector and the ground station receiver were
found to be using different topic names. Both processes ran without error,
and neither reported that it had no counterpart; the mismatch was visible
only as a publisher with zero subscribers and a subscriber with zero
publishers. Counting matched endpoints is consequently treated as a required
test rather than a diagnostic step.

\subsubsection{The send-buffer fault: a worked example}
\label{sec:test-wmem}

Test~T3.2 is the only test in this campaign recorded as having failed in the
delivered system, and it is documented in full because it demonstrates the
argument of Section~\ref{sec:test-strategy} more completely than any other
fault encountered.

The symptom was that the edge node received imagery at \(0.19\)\,Hz where the
camera was configured for \(5\)\,Hz --- a shortfall of approximately
\(26\times\). Every plausible explanation was eliminated in turn, and each
elimination was itself a measurement.

\begin{table}[htbp]
\centering
\caption{Eliminating candidate causes of the frame-rate shortfall. Each
measurement excludes one layer.}
\label{tab:test-wmem}
\begin{tabular}{@{}p{34mm}p{40mm}p{40mm}@{}}
\toprule
Hypothesis & Measurement & Outcome \\
\midrule
Simulation running slowly
  & Real-time factor from \texttt{/clock}
  & \(0.998\) --- simulated time tracks wall clock; excluded \\[3pt]
Gazebo rendering too slowly
  & Host-side subscriber on the camera topic
  & \(5.00\)\,Hz, gaps \(0.19\)--\(0.21\)\,s; source is correct \\[3pt]
Software rendering
  & GL driver mapped by \texttt{gzserver}
  & Hardware driver loaded; excluded \\[3pt]
Link negotiated at $100$\,Mbps
  & \texttt{ethtool} on the edge node
  & $1000$\,Mb/s full duplex; excluded \\[3pt]
Fragments lost in flight
  & Interface counters on both machines
  & Zero errors, zero drops on both; excluded \\[3pt]
Receive buffers too small
  & \texttt{net.core.rmem\_max}
  & $512$\,MB, already raised; excluded \\[3pt]
\textbf{Send buffers too small}
  & \texttt{net.core.wmem\_max}
  & \(\mathbf{208}\)\,\textbf{kB} --- default, never raised \\
\bottomrule
\end{tabular}
\end{table}

The DDS profile requests a \(16\)\,MB send buffer. The kernel caps any such
request at \texttt{net.core.wmem\_max}, which stood at its \(208\)\,kB
default, and performs the reduction silently. A single frame is therefore
approximately \(13\times\) larger than the buffer available to transmit it.
The publisher filled that buffer immediately and discarded the remaining
fragments at the socket layer, before they reached the network interface.

Confirmation came from comparing what was produced with what was
transmitted. The host generated \(414\)\,MB of imagery over thirty seconds
and placed only \(10\)\,MB of it on the wire --- \(2\)\,Mbps against the
\(110\)\,Mbps the configuration implies. Raising
\texttt{net.core.wmem\_max} to \(512\)\,MB on both machines and restarting
the stack, which is necessary because Fast DDS applies socket options only
at participant creation, restored delivery to \(5.000\)\,Hz with a standard
deviation of \(4.5\)\,ms.

Three properties of this fault make it worth recording.

\paragraph{No counter records it.}
Loss at the socket layer is not visible to interface statistics.
\texttt{tx\_dropped} on the sending interface and \texttt{rx\_errors} on the
receiving one both remained at zero throughout, because the discarded
fragments were never transmitted. A diagnosis founded on interface counters
alone would conclude that the network was healthy, which it was.

\paragraph{Its symptom appeared three layers above its cause.}
The visible consequences were that the detector reported no humans in
\(93\) of \(94\) frames and that the patrol mission timed out at every
waypoint. Both invite explanation in terms of detector configuration or
flight control. Neither has anything to do with either: the aircraft flew
correctly and the detector functioned correctly, but the aircraft was
observed through roughly one frame every five seconds.

\paragraph{Asymmetric configuration is easy to overlook.}
The receive path had been corrected during earlier work and the transmit
path had not, because the failure that prompted the original correction ---
incomplete reassembly at the subscriber --- pointed only at the receiving
side. The two directions are configured by separate kernel parameters, and
an error in either produces the same observable outcome. The pre-flight
checks in \texttt{scripts/netns/run\_hitl.sh} were subsequently extended to
verify both parameters on both machines.

\begin{figure}[htbp]
\centering
\includegraphics[width=0.85\textwidth]{camera_rate_before_after.png}
\caption{Delivered camera rate at the edge node, measured with
\texttt{ros2 topic hz}, before and after \texttt{net.core.wmem\_max} was
raised. The source rate was \(5\)\,Hz throughout and no other parameter was
changed. Delivery rose from \(0.19\)\,Hz to \(5.000\)\,Hz, and the
inter-arrival standard deviation fell from over a second to \(4.5\)\,ms.}
\label{fig:test-ratefix}
\end{figure}

% ---------------------------------------------------------------------
\subsection{Layer 4 --- Detector function}
\label{sec:test-l4}

This layer tests that the detector produces correct output, independently of
how that output is transported. It is the only layer with an external ground
truth: the number and position of the human models are fixed in the world
definition, so recall is a directly measurable quantity rather than an
estimate.

\begin{table}[htbp]
\centering
\caption{Layer 4 test cases --- detector function.}
\label{tab:test-l4}
\begin{tabular}{@{}lp{40mm}p{38mm}l@{}}
\toprule
ID & Test and command & Result & Verdict \\
\midrule
T4.1 & Model loads and runs on ARM. \newline
       Run detector standalone on a saved image
     & Inference completes, $623$\,ms
     & \pass \\[2pt]
T4.2 & Detections are on real subjects. \newline
       View \texttt{/uav1/camera/annotated}
     & Bounding boxes coincide with the human models
     & \pass \\[2pt]
T4.3 & False positive rate. \newline
       Count detections against ground truth of $5$
     & Earlier configuration reported up to $16$; corrected configuration
       reports $2$--$4$
     & \pass \\[2pt]
T4.4 & Recall against known ground truth. \newline
       Count detections against ground truth of $5$
     & $2$--$4$ of $5$, i.e.\ $40$--$80\%$
     & \pass \\[2pt]
T4.5 & Message content and size. \newline
       \texttt{ros2 topic echo /detections/uav1}
     & JSON, approximately $118$\,B, human count present
     & \pass \\
\bottomrule
\end{tabular}
\end{table}

\begin{figure}[htbp]
\centering
\begin{subfigure}[t]{0.48\textwidth}
  \centering
  \includegraphics[width=\textwidth]{world_ground_truth.png}
  \caption{The test scene in Gazebo. Five human models stand approximately
  one metre apart, at coordinates fixed in the world definition.}
  \label{fig:test-groundtruth}
\end{subfigure}
\hfill
\begin{subfigure}[t]{0.48\textwidth}
  \centering
  \includegraphics[width=\textwidth]{detection_annotated.png}
  \caption{Detector output. Bounding boxes coincide with the models,
  confirming that reported counts correspond to real subjects.}
  \label{fig:test-annotated}
\end{subfigure}
\caption{Ground truth and detector output. Because the number and position
of the subjects are fixed in the world definition, recall is counted rather
than estimated, and any count exceeding five is necessarily a false
positive.}
\label{fig:test-detection}
\end{figure}

Test~T4.2 was made a precondition of accepting any quantitative result. A
detector that reports a plausible count while placing its boxes on building
texture would satisfy every numerical check in this section, and the earlier
configuration recorded in T4.3 did exactly that: sixteen detections against
five subjects, all confident, none correct. Visual inspection of annotated
frames, shown in Figure~\ref{fig:test-detection}, is the only test that
distinguishes the two cases, and it was therefore run before the timing
experiments rather than after them. The false positives were eliminated by
raising the confidence threshold and narrowing the camera field of view.

The shortfall in T4.4 is attributed to scene geometry rather than to the
detector. The five models stand approximately one metre apart, so from the
oblique camera view the nearer figures partially occlude those behind them,
and a partially occluded subject at approximately $53$ pixels in height
falls below the confidence threshold. This is a property of the test scene,
and is noted as such rather than presented as a detector accuracy figure ---
the project measures the cost of onboard detection, not the accuracy of
YOLOv8n, which is established elsewhere.

% ---------------------------------------------------------------------
\subsection{Layer 5 --- End-to-end path}
\label{sec:test-l5}

The final layer tests the complete path: Gazebo renders a frame, the frame
crosses the unimpaired wired link to the Pi, the Pi performs detection, and
the result crosses the simulated wireless channel to the ground station.
This is the behaviour the system exists to demonstrate, and it is the only
layer at which the requirement --- that results, and only results, cross the
radio --- can be confirmed as a whole.

\begin{table}[htbp]
\centering
\caption{Layer 5 test cases --- end-to-end operation.}
\label{tab:test-l5}
\begin{tabular}{@{}lp{40mm}p{38mm}l@{}}
\toprule
ID & Test and command & Result & Verdict \\
\midrule
T5.1 & Whole stack starts from one command. \newline
       \texttt{run\_hitl.sh -{}-mission}
     & All stages report ready; mission flies
     & \pass \\[2pt]
T5.2 & Detections arrive at the ground station. \newline
       Ground station receiver log, inside \texttt{gcsns}
     & Human counts received during flight
     & \pass \\[2pt]
T5.3 & Those detections crossed ns-3. \newline
       Receiver runs in \texttt{gcsns}, reachable only via
       \texttt{tap-gcs}
     & Namespace has no other route to the Pi
     & \pass \\[2pt]
T5.4 & Imagery does not cross ns-3. \newline
       As T2.5, T2.6
     & Only detection-sized traffic on \texttt{tap-uav2}
     & \pass \\[2pt]
T5.5 & Clean shutdown. \newline
       \texttt{Ctrl+C}, then \texttt{ip -br link}
     & Remote detector stopped; namespaces, bridges and VLANs removed
     & \pass \\
\bottomrule
\end{tabular}
\end{table}

Test~T5.3 is the structural argument that supports the claim made throughout
this report. The ground station receiver executes inside the \texttt{gcsns}
network namespace, whose only interface is one end of a veth pair attached
to \texttt{br-gcs} and thence to \texttt{tap-gcs}. The namespace has no
route to $10.0.0.0/24$ and no interface on it. A message that arrives at
that receiver cannot have taken any path other than through ns-3, because no
other path exists from inside the namespace. Separation is therefore a
property of the topology rather than of a configuration setting that could
be silently overridden, which is a stronger guarantee than filtering would
provide.

% ---------------------------------------------------------------------
\subsection{Experimentation}
\label{sec:test-experiments}

Three controlled experiments were performed once all five layers passed. In
each, one parameter was varied and the remainder held fixed.

\paragraph{E1 --- Inference back-end and input resolution.}
The detector was timed across the four combinations of back-end (PyTorch,
NCNN) and input resolution ($960 \times 544$, $640 \times 384$), reported in
Table~\ref{tab:hitl-inference}. Changing back-end at constant resolution
gives approximately $2.5\times$; halving the pixel count at constant
back-end gives approximately $2.3\times$. Latency therefore scales close to
linearly with pixel count, a pixel ratio of $0.47$ producing a time ratio of
$0.43$. Timing was taken around the model invocation alone so that the two
back-ends are compared on identical work.

\paragraph{E2 --- Source frame rate.}
The camera rate was reduced at source from $20$\,Hz to $5$\,Hz, still
approximately seven times the rate the detector can consume, so that no
detection opportunity is lost. Inference latency fell from $1{,}343$\,ms to
approximately $1{,}000$\,ms and link occupancy from $481$\,Mbps to
$111$\,Mbps. The independent variable is deliberately the source rate rather
than a rate limit applied on arrival: discarding frames after reception
leaves the transport and deserialisation cost already incurred, and so does
not test the same thing.

\paragraph{E3 --- Isolated versus in-pipeline inference.}
The same model, on the same input, was timed standalone on a saved image and
within the live pipeline: $623$\,ms against $1{,}343$\,ms. The $720$\,ms
difference is the interference between the two workloads, and amounts to
approximately half the board's available inference throughput. This is the
principal finding of the hardware-in-the-loop work, and it is reported in
full in Section~\ref{sec:hitl-contention}.

% ---------------------------------------------------------------------
\subsection{Validation of the measurements}
\label{sec:test-validation}

Passing a test establishes that a component behaves as specified. It does
not establish that a measurement taken from that component means what it is
claimed to mean. Four validation arguments are therefore made explicitly.

\paragraph{Sample size.}
Latency figures are reported over thirty ICMP samples. A two-sample
measurement of the same configuration returned a mean of $167$\,ms against
$59$\,ms for thirty, because the first packet includes address resolution
across the simulated channel. The short sample is not noisy but biased, and
consistently so; increasing the sample size is the correct remedy rather
than repetition.

\paragraph{Placement of the timer.}
Inference duration is measured around the model invocation alone, excluding
message reception, format conversion and publication. This placement is what
makes E3 interpretable. Had the timer enclosed the whole callback, the
$720$\,ms difference could equally have been queuing delay ahead of the
model. Because it does not, the additional time is spent inside the model
call itself, which establishes contention for processor time rather than
waiting.

\paragraph{Ground truth.}
Detection performance is validated against the world definition, in which
the number and position of the human models are fixed and known. Recall is
therefore counted, not estimated. The same fact supports T4.3: a count
exceeding five is necessarily a false positive, without any need to
adjudicate individual boxes.

\paragraph{Rate measurements presume a running clock.}
Gazebo runs in lockstep with the autopilot, advancing only when SITL sends
servo commands. If SITL stalls, simulated time stops, and every rate derived
from it becomes meaningless while remaining superficially plausible. During
one diagnosis a low frame rate was attributed in turn to software rendering
and to camera geometry before a stalled clock was found to be the cause.
The state of \texttt{/clock} is consequently checked before any rate figure
is accepted, and this is recorded here as a methodological requirement of
lockstep simulation rather than as an incidental fault.

% ---------------------------------------------------------------------
\subsection{Coverage and outstanding items}
\label{sec:test-coverage}

Table~\ref{tab:test-coverage} summarises coverage by layer. Every layer
required for the reported results has been tested and passed.

\begin{table}[htbp]
\centering
\caption{Test coverage by layer.}
\label{tab:test-coverage}
\begin{tabular}{@{}llcl@{}}
\toprule
Layer & Scope & Cases & Status \\
\midrule
L1 & Physical link            & 4 & All passing \\
L2 & Path separation          & 6 & All passing \\
L3 & Middleware transport     & 7 & All passing after T3.2 corrected \\
L4 & Detector function        & 5 & All passing \\
L5 & End-to-end operation     & 5 & All passing \\
\bottomrule
\end{tabular}
\end{table}

Three items remain outside the tested set, and no result in this report
depends on them.

Clock synchronisation between the two machines is not established. End-to-end
latency would be computed from timestamps taken on different hosts, and
without synchronisation such values are not meaningful while remaining
plausible in magnitude. End-to-end latency is therefore not reported, rather
than reported with a caveat.

The ground configuration --- compressed imagery crossing the simulated
channel with inference performed at the ground station --- has not been
measured. The comparison it enables is the subject of the remaining work.

Behaviour at range has not been characterised. All measurements were taken
with the aircraft close to the ground station, giving zero packet loss on
both paths. The reliability policy applied to detection messages is
consequently untested under loss, which is the condition it exists to
handle.

% =====================================================================
%  IMAGE CHECKLIST
%  Save each file into  docs/report/figures/  under the name given.
%
%  1. testbed_hardware.jpg          (Fig. \ref{fig:test-hardware})
%     A photograph of the real setup: laptop, USB gigabit adapter, the
%     Ethernet cable, and the Raspberry Pi 4B with its power supply.
%     Shoot it so the single cable between the two machines is clearly
%     the only link. This is the one image that shows the work is
%     hardware-in-the-loop rather than pure simulation, so it is worth
%     taking properly: good light, plain background, both machines in
%     frame. Optionally add small printed labels reading "VLAN 10 /
%     VLAN 42" beside the cable.
%
%  2. world_ground_truth.png        (Fig. \ref{fig:test-groundtruth})
%     Gazebo client view of the five human models, from above or at an
%     angle that makes all five countable. Purpose: the reader must be
%     able to verify that the ground truth really is five.
%         gzclient        # with the stack running
%
%  3. detection_annotated.png       (Fig. \ref{fig:test-annotated})
%     A frame with bounding boxes drawn on real people. Since the
%     annotated topic was removed from the edge node, capture this by
%     running the detector once with show_window enabled, or by drawing
%     the boxes host-side from the bbox values in /detections/uav1.
%     Choose a frame where the boxes clearly sit on the models.
%
%  4. camera_rate_before_after.png  (Fig. \ref{fig:test-ratefix})
%     Evidence for the send-buffer fault. Either a screenshot of the two
%     `ros2 topic hz` outputs side by side (0.19 Hz and 5.000 Hz), or a
%     small plot of inter-arrival time against sample index for both
%     runs. The plot is stronger: it shows the variance collapsing as
%     well as the rate rising.
%
%  OPTIONAL, if space allows
%  5. A screenshot of the two ping results (10.0.0.1 and 10.42.0.10)
%     side by side. Table \ref{tab:test-l2} already carries the numbers,
%     so this is supporting evidence rather than a result.
%  6. rqt_graph of the running node/topic graph, which shows at a glance
%     that the detector and the ground station are connected only
%     through /detections/uav1.
% =====================================================================
